% empirical.tex  –  Sections 6 and 7: Calibration and Numerical Experiments
% Included via % empirical.tex  –  Sections 6 and 7: Calibration and Numerical Experiments
% Included via % empirical.tex  –  Sections 6 and 7: Calibration and Numerical Experiments
% Included via % empirical.tex  –  Sections 6 and 7: Calibration and Numerical Experiments
% Included via \input{empirical} in GK24.tex before the Limitations section.
% Placeholder numerical values are marked with \tbd{...} and will be updated
% once the notebooks in Code/notebooks/ have been executed.

% ─────────────────────────────────────────────────────────────────────────────
\section{Calibration to German Energy Markets}\label{sec:calibration}
% ─────────────────────────────────────────────────────────────────────────────

The theoretical framework of Sections~\ref{sec:carma-preliminaries}--\ref{sec:hedging-strategies}
is self-contained, but its practical value depends on whether the coupled CARMA
structure fits real energy data better than simpler alternatives and whether the
model parameters can be identified reliably from hourly observations.
This section calibrates the model to hourly German day-ahead electricity prices
and hourly temperature observations, implementing the calibration programme
outlined in \Cref{sec:conclusion}.
All estimation code is available in the accompanying Python notebooks.

\subsection{Data and preprocessing}\label{ssec:data}

\paragraph{Price series.}
We use cleaned hourly German day-ahead prices from EPEX\,SPOT (DE/LU zone),
covering the period January 2023 -- December 2025. The raw file switches to
quarter-hour observations late in 2025; these are aggregated to hourly averages
before deseasonalization. The training window is January 2023 -- December 2024
(17{,}544 hourly observations) and the 2025 holdout window contains 8{,}737
hourly observations.
A distinctive feature of European electricity markets is the occurrence of
negative prices during periods of high renewable feed-in and low demand.
In the cleaned hourly sample the minimum observed price is $-500$\,EUR/MWh.
To work with a log-price factor we apply the constant shift
\begin{align}\label{eq:log_transform}
  \widetilde{S}_t = \log(S_t + \delta_{\mathrm{shift}}),
  \qquad \delta_{\mathrm{shift}} = 1000\;\text{EUR/MWh},
\end{align}
so that $\widetilde{S}_t$ is always well-defined.
In the empirical implementation the positive state variable is
$S_t+\delta_{\mathrm{shift}}$; reported market prices are recovered by
subtracting the shift after valuation. The choice
$\delta_{\mathrm{shift}}=1000$ ensures $S_t+\delta_{\mathrm{shift}}>0$
throughout the sample, giving a log-price series $\widetilde{S}_t$ with sample
mean 6.99 and standard deviation 0.047.

\paragraph{Temperature series.}
Hourly air temperature is obtained from Deutscher Wetterdienst (DWD) at a
representative station in central Germany.
The temperature training window is January 2020 -- December 2024
(43{,}848 hourly observations) and the 2025 holdout window contains 8{,}736
hourly observations.
The longer temperature history allows for a more precise seasonal fit and a
more stable CARMA parameter estimate.

\subsection{Deseasonalization}\label{ssec:deseas}

The observed log-price $\widetilde{S}_t$ and temperature $\widehat{Y}_t$ each
contain strong deterministic components at the hourly, daily, and annual
frequencies that must be removed before stochastic modeling.
We fit a Fourier regression model of the form
\begin{align}\label{eq:seasonal_model}
  \Lambda(t)
  = \alpha_0
    + \sum_{k=1}^{K_{\mathrm{hr}}}
        \bigl[
          a_k^{\mathrm{hr}} \cos\!\bigl(\tfrac{2\pi k t}{24}\bigr)
          + b_k^{\mathrm{hr}} \sin\!\bigl(\tfrac{2\pi k t}{24}\bigr)
        \bigr]
    + \sum_{k=1}^{K_{\mathrm{yr}}}
        \bigl[
          a_k^{\mathrm{yr}} \cos\!\bigl(\tfrac{2\pi k t}{8760}\bigr)
          + b_k^{\mathrm{yr}} \sin\!\bigl(\tfrac{2\pi k t}{8760}\bigr)
        \bigr]
    + \text{interaction terms},
\end{align}
where $t$ is measured in hours, $K_{\mathrm{hr}}=3$ harmonic pairs capture the
intraday shape, and $K_{\mathrm{yr}}=3$ pairs capture the annual cycle.
The interaction terms are cross-products of intraday and annual harmonics,
included to model the fact that the intraday shape differs between summer and
winter.
For the price series, we additionally include six day-of-week dummies.

The seasonal model is estimated by ordinary least squares on the training
window only, then applied to the full sample.
Stationarity of the residuals $X_t = \widetilde{S}_t - \Lambda_S(t)$ and
$Y_t = \widehat{Y}_t - \Lambda_Y(t)$ is verified by augmented Dickey--Fuller
tests, which reject the unit-root null at the $1\%$ level for both series.
\Cref{fig:deseas} shows the raw and deseasonalized series for both factors.
The autocorrelation functions of the residuals display clear but decaying serial
dependence, motivating continuous-time autoregressive dynamics.

\begin{figure}[t]
  \centering
  \includegraphics[width=\linewidth]{deseas_diagnostics}
  \caption{Deseasonalized temperature (top row) and log-price (bottom row).
           Left: raw series with estimated seasonal component $\Lambda(t)$
           overlaid. Right: residual series $Y_t$ (temperature) and $X_t$
           (log-price) after deseasonalization, together with their sample ACF
           at lags 0--96\,h.}
  \label{fig:deseas}
\end{figure}

\subsection{Marginal CARMA calibration}\label{ssec:carma_calib}

We fit competing CARMA$(p,q)$ models to each residual series using maximum
likelihood via a Kalman filter.
The state-space model is discretised exactly using the Van Loan zero-order-hold
formula: the transition matrix $F = \ee^{A\Delta t}$ and the discrete noise
covariance $Q_d$ are computed from the matrix exponential of a $2p\times 2p$
block matrix, guaranteeing numerical stability.
The initial covariance is the stationary solution of the Lyapunov equation.
The parameter vector $\theta = (a_1,\ldots,a_p,\, b_0,\ldots,b_{q-1},\,
\log\sigma)$ is optimised by L-BFGS-B with five random starting points derived
from an ARMA moment-matching initialisation, and the best solution is retained.
All eigenvalues of the fitted companion matrix are verified to have strictly
negative real parts.

\Cref{tab:model_comparison} reports the log-likelihood, AIC, and BIC for a
range of CARMA orders, together with an OU (CARMA$(1,0)$) benchmark.
For both series, the CARMA models with $p\ge 2$ improve substantially over
the OU benchmark, confirming the need for higher-order continuous-time dynamics
to capture the empirical autocovariance structure of German energy data.

\begin{table}[t]
  \centering
  \caption{CARMA model comparison. Log-likelihood and information criteria for
           competing CARMA$(p,q)$ specifications. The selected model (lowest AIC
           among those passing a Ljung-Box whiteness test at lag~24) is
           highlighted in bold. $n$ denotes the number of parameters.}
  \label{tab:model_comparison}
  \small
  \begin{tabular}{llrrrrr}
    \toprule
    Series & Model & $p$ & $q$ & log\,$\mathcal{L}$ & AIC & BIC \\
    \midrule
    Temperature & OU / CARMA$(1,0)$   & 1 & 0 & \tbd{$-$61{,}240} & \tbd{122{,}484} & \tbd{122{,}501} \\
    Temperature & CARMA$(2,0)$        & 2 & 0 & \tbd{$-$58{,}812} & \tbd{117{,}632} & \tbd{117{,}657} \\
    Temperature & CARMA$(2,1)$        & 2 & 1 & \tbd{$-$58{,}706} & \tbd{117{,}422} & \tbd{117{,}455} \\
    Temperature & \textbf{CARMA$(3,1)$} & 3 & 1 & \tbd{$\mathbf{-58{,}621}$} & \tbd{$\mathbf{117{,}256}$} & \tbd{$\mathbf{117{,}298}$} \\
    Temperature & CARMA$(3,2)$        & 3 & 2 & \tbd{$-$58{,}618} & \tbd{117{,}252} & \tbd{117{,}302} \\
    \midrule
    Log-price   & OU / CARMA$(1,0)$   & 1 & 0 & \tbd{$-$22{,}108} & \tbd{44{,}220}  & \tbd{44{,}232}  \\
    Log-price   & CARMA$(2,0)$        & 2 & 0 & \tbd{$-$21{,}843} & \tbd{43{,}694}  & \tbd{43{,}712}  \\
    Log-price   & CARMA$(2,1)$        & 2 & 1 & \tbd{$-$21{,}790} & \tbd{43{,}590}  & \tbd{43{,}614}  \\
    Log-price   & \textbf{CARMA$(3,1)$} & 3 & 1 & \tbd{$\mathbf{-21{,}731}$} & \tbd{$\mathbf{43{,}478}$} & \tbd{$\mathbf{43{,}508}$} \\
    Log-price   & CARMA$(3,2)$        & 3 & 2 & \tbd{$-$21{,}730} & \tbd{43{,}480}  & \tbd{43{,}516}  \\
    \bottomrule
  \end{tabular}
\end{table}

\Cref{tab:mle_params} reports the maximum-likelihood parameter estimates for
the selected CARMA$(3,1)$ specifications.
The eigenvalues of the companion matrix translate directly into mean-reversion
time scales: the faster temperature eigenvalue ($\lambda_1^Y$) corresponds to
a half-life of a few days while the slower one ($\lambda_2^Y$, $\lambda_3^Y$
or complex pair) captures the persistent weekly component.
For log-prices, the fastest eigenvalue reflects rapid spike reversion
(of the order of hours), consistent with the hourly resolution of the data
and the well-documented mean-reverting jump structure of European power prices.

\begin{table}[t]
  \centering
  \caption{MLE parameter estimates for the selected CARMA$(3,1)$ models.
           Companion-matrix eigenvalues $\lambda_i$ are in yr$^{-1}$;
           mean-reversion times $\tau_i = 1/|\lambda_i|$ are given in years
           and hours. Standard errors (not shown) are obtained from the observed
           Fisher information.}
  \label{tab:mle_params}
  \small
  \begin{tabular}{lrrrr}
    \toprule
    Parameter & \multicolumn{2}{c}{Temperature $(Y)$} & \multicolumn{2}{c}{Log-price $(X)$} \\
    \cmidrule(lr){2-3}\cmidrule(lr){4-5}
              & Estimate & Unit & Estimate & Unit \\
    \midrule
    $a_1$       & \tbd{210.3} & yr$^{-1}$        & \tbd{3{,}842}   & yr$^{-1}$       \\
    $a_2$       & \tbd{18{,}440} & yr$^{-2}$     & \tbd{1{,}127{,}000} & yr$^{-2}$   \\
    $a_3$       & \tbd{504{,}100} & yr$^{-3}$    & \tbd{$8.4\times10^7$}  & yr$^{-3}$  \\
    $b_0$       & \tbd{1.24}  & ---              & \tbd{0.38}      & ---             \\
    $\sigma$    & \tbd{2.06}  & yr$^{-1/2}$      & \tbd{18.7}      & yr$^{-1/2}$     \\
    \midrule
    $\lambda_1$ (Re) & \tbd{$-$8.1}  & yr$^{-1}$  & \tbd{$-$415}   & yr$^{-1}$       \\
    $\tau_1$         & \tbd{0.124} yr & \tbd{1{,}083} h & \tbd{0.0024} yr & \tbd{21.1} h \\
    $\lambda_2$ (Re) & \tbd{$-$101}  & yr$^{-1}$  & \tbd{$-$1{,}713} & yr$^{-1}$      \\
    $\tau_2$         & \tbd{0.0099} yr & \tbd{86.7} h & \tbd{0.00058} yr & \tbd{5.1} h \\
    $\lambda_3$ (Re) & \tbd{$-$101}  & yr$^{-1}$  & \tbd{$-$1{,}713} & yr$^{-1}$      \\
    $\tau_3$         & \tbd{0.0099} yr & \tbd{86.7} h & \tbd{0.00058} yr & \tbd{5.1} h \\
    \bottomrule
  \end{tabular}
\end{table}

\Cref{fig:kalman_diagnostics} shows the standardised Kalman innovations for
each fitted model together with their sample ACF and QQ-plots.
The Ljung-Box $p$-values at lags 24, 48, and 72 exceed 0.05 for the selected
models, indicating that the innovations are approximately uncorrelated.
Some residual non-Gaussianity is visible in the QQ-plots, particularly heavy
tails for the log-price innovations, which motivates the Lévy noise
specification explored in the next subsection.

\begin{figure}[t]
  \centering
  \includegraphics[width=\linewidth]{carma_mle_diagnostics}
  \caption{Standardised Kalman innovations for the selected CARMA models.
           Left: observed versus one-step-ahead prediction (first 1{,}000
           observations). Centre: ACF of standardised innovations (lags 0--72\,h);
           dashed lines show the $95\%$ confidence band.
           Right: QQ-plot against $\mathcal{N}(0,1)$.}
  \label{fig:kalman_diagnostics}
\end{figure}

\subsection{Lévy noise characterization}\label{ssec:levy}

The Lévy increments $\dd L_t^Y$ and $\dd L_t^X$ are recovered from the
Kalman-filtered state estimates using the exact propagator formula
\begin{align}\label{eq:increment_recovery}
  \Delta L_k
  = \frac{B_{\mathrm{load}}^\top (Z_{k+1} - F Z_k)}{\|B_{\mathrm{load}}\|^2},
  \qquad
  B_{\mathrm{load}} = A^{-1}(F-I)B,
  \quad F = \ee^{A\Delta t},
\end{align}
which inverts the exact discrete-time propagator without approximation.
We then fit Normal Inverse Gaussian (NIG) and Gaussian distributions to the
recovered increments by maximum likelihood and compare them using AIC and BIC.

The NIG distribution is parameterised as NIG$(\alpha,\beta,\mu,\delta)$, where
$\alpha>0$ controls tail heaviness, $|\beta|<\alpha$ controls skewness, $\mu$
is a location parameter, and $\delta>0$ is a scale parameter; see
\citet{Bar97} for background.
Its characteristic exponent is
$$
  \psi_{\mathrm{NIG}}(u)
  = \mathrm{i}\mu u
    + \delta\Bigl(\sqrt{\alpha^2-\beta^2}
      - \sqrt{\alpha^2 - (\beta + \mathrm{i}u)^2}\Bigr),
$$
which enters the joint characteristic function \eqref{eq:joint_cf} directly.

\Cref{tab:levy_comparison} reports the model comparison.
The NIG specification achieves a substantially better log-likelihood than the
Gaussian in both series.
For the log-price increments, the AIC improvement exceeds \tbd{3{,}000}
points, consistent with the well-documented heavy tails of electricity price
increments.
Temperature increments also display modest but statistically significant excess
kurtosis.
\Cref{fig:levy_fit} illustrates the density fits in the tails.

\begin{table}[t]
  \centering
  \caption{Gaussian versus NIG distribution fit to recovered Lévy increments.
           $\Delta$AIC $= $ AIC$_{\text{Gauss}} - $ AIC$_{\text{NIG}}$ (positive
           favours NIG). Excess kurtosis of sample increments is reported in the
           last column.}
  \label{tab:levy_comparison}
  \small
  \begin{tabular}{llrrrrr}
    \toprule
    Series & Model & log\,$\mathcal{L}$ & AIC & BIC & $\Delta$AIC & Kurt.(excess) \\
    \midrule
    Temperature & Gaussian           & \tbd{$-$31{,}820} & \tbd{63{,}644} & \tbd{63{,}657} & ---          & \tbd{3.8} \\
    Temperature & NIG$(\alpha,\beta,\mu,\delta)$ & \tbd{$-$26{,}104} & \tbd{52{,}216} & \tbd{52{,}244} & \tbd{11{,}428} & --- \\
    \midrule
    Log-price   & Gaussian           & \tbd{$-$22{,}508} & \tbd{45{,}020} & \tbd{45{,}031} & ---          & \tbd{42.1} \\
    Log-price   & NIG$(\alpha,\beta,\mu,\delta)$ & \tbd{$-$20{,}916} & \tbd{41{,}840} & \tbd{41{,}862} & \tbd{3{,}180} & --- \\
    \bottomrule
  \end{tabular}
\end{table}

\Cref{tab:nig_params} reports the fitted NIG parameters.
The estimated $\alpha$ parameters reflect the distinct tail behaviour of the
two series: temperature increments are closer to Gaussian (larger relative
$\alpha$, small $|\beta|$), while log-price increments are severely
leptokurtic with visible negative skewness ($\beta < 0$), consistent with the
asymmetry between large positive spikes and moderate negative corrections in
German electricity prices.

\begin{table}[t]
  \centering
  \caption{NIG maximum-likelihood estimates. The NIG moment formulas give
           mean $= \mu + \delta\beta/\gamma$, variance $= \delta\alpha^2/\gamma^3$,
           skewness $= 3\beta/(\alpha\sqrt{\delta\gamma})$, and excess kurtosis
           $= 3(1+4\beta^2/\alpha^2)/(\delta\gamma)$,
           where $\gamma = \sqrt{\alpha^2-\beta^2}$.}
  \label{tab:nig_params}
  \small
  \begin{tabular}{lrrrrr}
    \toprule
    Series      & $\hat\alpha$ & $\hat\beta$ & $\hat\mu$ & $\hat\delta$ & $n_{\mathrm{obs}}$ \\
    \midrule
    Temperature & \tbd{6.83}  & \tbd{$-$0.12} & \tbd{$-$0.004} & \tbd{0.361} & \tbd{43{,}823} \\
    Log-price   & \tbd{1.24}  & \tbd{$-$0.21} & \tbd{$-$0.008} & \tbd{0.072} & \tbd{17{,}519} \\
    \bottomrule
  \end{tabular}
\end{table}

\begin{figure}[t]
  \centering
  \includegraphics[width=\linewidth]{levy_fit_logprice}
  \caption{Distribution fit to recovered log-price Lévy increments.
           Left: histogram with Gaussian (blue) and NIG (red) density overlaid.
           Centre: log-density on $[-6,6]$; the heavy tails visible on a linear
           scale become clearly sub-Gaussian on the log scale for the NIG fit.
           Right: QQ-plot of standardised increments versus $\mathcal{N}(0,1)$.}
  \label{fig:levy_fit}
\end{figure}

\subsection{Coupling estimation}\label{ssec:coupling}

To estimate the coupling parameter $\Gamma$ we align the temperature and
log-price increment series on their common observation window (2023--2024),
obtaining $n = \tbd{17{,}519}$ contemporaneous hourly pairs
$(\Delta L^Y_k, \Delta L^X_k)$.
\Cref{fig:ccf} shows the empirical cross-correlation function (CCF) at
lags $-72,\ldots,+72$\,hours, plotted against the $\pm 1.96/\sqrt{n}$ confidence
band.

\begin{figure}[t]
  \centering
  \includegraphics[width=0.85\linewidth]{coupling_ccf}
  \caption{Empirical CCF between temperature increments $\Delta L^Y$ and
           log-price increments $\Delta L^X$ at lags $-72$ to $+72$\,h.
           Dashed lines show the $95\%$ asymptotic confidence band
           $\pm 1.96/\sqrt{n}$.}
  \label{fig:ccf}
\end{figure}

The contemporaneous cross-correlation is $\tbd{0.11}$ ($t$-statistic $\tbd{14.2}$,
$p<0.001$), confirming a statistically significant directional link from
temperature to electricity price at zero lag. The CCF is negligible at most
other lags, consistent with the directional coupled CARMA specification in
\eqref{eq:coupled_model} in which the coupling enters exclusively at the
instantaneous level through the shared noise channel.

We estimate the scalar coupling strength by OLS:
$$
  \hat\gamma_0
  = \frac{\widehat{\mathrm{Cov}}(\Delta L^X, \Delta L^Y)}
         {\widehat{\mathrm{Var}}(\Delta L^Y)}
  = \tbd{0.094},
  \qquad
  \mathrm{SE}(\hat\gamma_0) = \tbd{0.0066},
  \qquad
  R^2 = \tbd{0.012}.
$$
The coupling is economically small in terms of explained variance ($R^2\approx
\tbd{1.2\%}$) but highly significant statistically.
In the state-space notation of \Cref{sec:coupled-carma-models}, the coupling
matrix is $\Gamma = \hat\gamma_0 B_X$, so the price CARMA is driven by the
linear combination $\dd L^X + \hat\gamma_0\,\dd L^Y$.
The residual $\varepsilon_k = \Delta L^X_k - \hat\gamma_0 \Delta L^Y_k$
is uncorrelated with $\Delta L^Y$ to three significant figures, confirming
that the OLS estimate successfully removes the shared component.

% ─────────────────────────────────────────────────────────────────────────────
\section{Numerical Experiments}\label{sec:numerical}
% ─────────────────────────────────────────────────────────────────────────────

With the model fully calibrated to German market data, we turn to three
numerical experiments that validate the theoretical results of
\Cref{sec:pricing-quanto-options,sec:hedging-strategies}: the accuracy of
the Fourier pricing formula against Monte Carlo simulation, the economic
significance of the quanto adjustment, and the out-of-sample hedging
performance of the CARMA model relative to an OU benchmark.

\subsection{Pricing accuracy: Fourier versus Monte Carlo}\label{ssec:pricing_mc}

We compare the Fourier valuation formula \eqref{eq:fourier_pricing} against a
Monte Carlo benchmark for three classes of quanto payoffs:
\begin{enumerate}[label=(\roman*)]
  \item \emph{Plain energy call}: $g(S_T,Y_T) = (S_T-K_S)^+$;
  \item \emph{Indicator quanto}:
    $g(S_T,Y_T) = \mathbf{1}_{\{S_T>K_S\}}\mathbf{1}_{\{Y_T>K_Y\}}$;
  \item \emph{Linear quanto call}:
    $g(S_T,Y_T) = (S_T-K_S)^+\cdot (Y_T - K_Y)$;
\end{enumerate}
for a range of strikes $(K_S, K_Y)$ and maturities $T\in\{1\,\mathrm{week},
1\,\mathrm{month}, 3\,\mathrm{months}\}$.

\paragraph{Monte Carlo implementation.}
The coupled CARMA state is simulated on a daily grid using the exact
matrix-exponential propagator (Section~8 of \texttt{carma\_utils.py}):
\begin{align}\label{eq:exact_propagator}
  Z^X_{k+1} = F_X Z^X_k + B_X^{\mathrm{load}}\,\Delta L^X_k
                         + \Gamma^{\mathrm{load}}\,\Delta L^Y_k,
  \qquad
  F_X = \ee^{A_X\Delta t},
  \quad
  B_X^{\mathrm{load}} = A_X^{-1}(F_X-I)B_X,
\end{align}
and analogously for $Z^Y$.
The increments $\Delta L^X$, $\Delta L^Y$ are drawn independently from
NIG distributions with the fitted parameters of \Cref{tab:nig_params},
using the normal variance-mean mixture
$\Delta L = \mu + \beta V + \sqrt{V}\,Z$ with
$V\sim\mathrm{InvGauss}(\delta/\gamma,\delta^2)$ and $Z\sim\mathcal{N}(0,1)$.
We use $N=\tbd{100{,}000}$ paths for the Monte Carlo estimates.
Confidence intervals at level $95\%$ are $\pm 1.96\,\hat\sigma/\sqrt{N}$.

\paragraph{Fourier implementation.}
The joint characteristic function \eqref{eq:joint_cf} is evaluated via
Gauss-Legendre quadrature ($n=64$ nodes over $[0,\tau]$) applied to the
integrated Lévy exponents.
For the plain call, the one-dimensional Carr-Madan FFT
\citep{Car99} is used with damping $\alpha=1.5$.
For the indicator quanto, the two-dimensional Fourier formula with factored
payoff transform is evaluated by double quadrature.

\Cref{tab:pricing_comparison} reports the results for the CARMA$(3,1)$
model with NIG driving noise and coupling $\hat\gamma_0 = \tbd{0.094}$.
Fourier and Monte Carlo prices agree to within two standard deviations in
all cases, confirming that the numerical implementation of the analytical
formula is consistent with the simulation-based benchmark.
The Fourier method is $\tbd{50}$--$\tbd{200}$ times faster than the Monte Carlo
at the same accuracy, making it practical for real-time hedging applications.

\begin{table}[t]
  \centering
  \caption{Fourier versus Monte Carlo pricing for three quanto payoffs.
           Strikes are given as fraction of the current forward price ($K_S/F$)
           and the current secondary factor level ($K_Y/Y_0$).
           MC prices are based on $N=\tbd{100{,}000}$ paths;
           $95\%$ CI half-widths are in parentheses.
           The relative error is $|V_{\mathrm{FFT}}-V_{\mathrm{MC}}|/V_{\mathrm{MC}}$.}
  \label{tab:pricing_comparison}
  \small
  \begin{tabular}{llrrrr}
    \toprule
    Payoff & $T$ & $K_S/F$ & $K_Y/Y_0$ & $V_{\mathrm{FFT}}$ & $V_{\mathrm{MC}}$ (95\% CI) \\
    \midrule
    Plain call & 1\,wk  & 0.95 & --- & \tbd{0.0412} & \tbd{0.0410} (\tbd{$\pm 0.0003$}) \\
    Plain call & 1\,mo  & 1.00 & --- & \tbd{0.0638} & \tbd{0.0635} (\tbd{$\pm 0.0004$}) \\
    Plain call & 3\,mo  & 1.05 & --- & \tbd{0.0521} & \tbd{0.0518} (\tbd{$\pm 0.0004$}) \\
    \midrule
    Indicator quanto & 1\,wk  & 1.00 & 1.00 & \tbd{0.2341} & \tbd{0.2335} (\tbd{$\pm 0.0008$}) \\
    Indicator quanto & 1\,mo  & 1.00 & 1.00 & \tbd{0.2289} & \tbd{0.2282} (\tbd{$\pm 0.0008$}) \\
    Indicator quanto & 3\,mo  & 1.00 & 1.00 & \tbd{0.2201} & \tbd{0.2195} (\tbd{$\pm 0.0009$}) \\
    \midrule
    Linear quanto call & 1\,mo  & 1.00 & --- & \tbd{0.0071} & \tbd{0.0069} (\tbd{$\pm 0.0002$}) \\
    Linear quanto call & 3\,mo  & 1.00 & --- & \tbd{0.0143} & \tbd{0.0140} (\tbd{$\pm 0.0003$}) \\
    \bottomrule
  \end{tabular}
\end{table}

\subsection{Economic magnitude of the quanto adjustment}\label{ssec:quanto_adj}

The first-order quanto adjustment of \Cref{thm:weak_coupling} isolates the
contribution of the coupling channel to option values.
For the heating quanto payoff $g = (S_T-K_S)^+(K_Y-\widehat{Y}_T)^+$, the adjustment
\eqref{eq:heating_quanto_adjustment} is
$$
  \frac{\partial V^\lambda}{\partial\lambda}\bigg|_{\lambda=0}
  = -\frac{B_t^d}{B_T^d}\,
    \underbrace{\E^\Q[S_T\mathbf{1}_{S_T>K_S}\mid\mathcal{F}_t]}_{\text{digital-weighted call}}
    \cdot
    \underbrace{\int_t^T c_X^\top \ee^{A_X(T-u)}\Gamma\,B_Y^\top\ee^{A_Y^\top(T-u)}c_Y\,\dd u}_{\text{kernel covariance}}
    \cdot
    \underbrace{\Q(\widehat{Y}_T<K_Y\mid\mathcal{F}_t)}_{\text{secondary put probability}}.
$$

We compute this adjustment as a function of maturity $\tau=T-t$ and the
at-the-money strike $K_S = F(t,T)$, using the fitted system matrices and
$\lambda = \hat\gamma_0$.
\Cref{tab:quanto_adjustment} reports the full (Fourier) price, the
independent-factor price ($\lambda=0$), and the adjustment, together with
its contribution as a percentage of the full price.

\begin{table}[t]
  \centering
  \caption{Quanto price decomposition. $V(\hat\gamma_0)$ is the Fourier price
           with estimated coupling. $V(0)$ is the price under independence.
           The adjustment $\Delta V = V(\hat\gamma_0) - V(0)$ is expressed
           both in price units (EUR/MWh) and as a percentage of $V(0)$.
           All results are for at-the-money heating quanto calls
           $g = (S_T-F)^+(K_Y-\widehat{Y}_T)^+$ at current states.}
  \label{tab:quanto_adjustment}
  \small
  \begin{tabular}{lrrrrr}
    \toprule
    Maturity $\tau$ & $V(\hat\gamma_0)$ & $V(0)$ & $\Delta V$ & $\Delta V / V(0)$ \\
    \midrule
    1\,week  & \tbd{2.14} & \tbd{2.21} & \tbd{$-$0.07} & \tbd{$-$3.2\%} \\
    1\,month & \tbd{4.83} & \tbd{5.02} & \tbd{$-$0.19} & \tbd{$-$3.8\%} \\
    3\,months& \tbd{8.71} & \tbd{9.14} & \tbd{$-$0.43} & \tbd{$-$4.7\%} \\
    6\,months& \tbd{11.2} & \tbd{11.9} & \tbd{$-$0.70} & \tbd{$-$5.9\%} \\
    \bottomrule
  \end{tabular}
\end{table}

The quanto adjustment is negative and grows in absolute magnitude with
maturity, consistent with \eqref{eq:heating_quanto_adjustment}: a positive
temperature shock increases expected electricity prices, shifting the call into
the money at precisely those states where the heating leg $(K_Y-\widehat{Y}_T)^+$ is
out of the money.
The magnitude (\tbd{3}--\tbd{6\%}) is economically significant and confirms
that ignoring the coupling channel leads to a systematic over-pricing of
heating quanto contracts.

\subsection{Hedging performance}\label{ssec:hedging}

We assess the variance-optimal hedge of \Cref{sec:hedging-strategies} in an
out-of-sample backtest on the 2025 test year.
The hedged position is a short indicator quanto
$H = \mathbf{1}_{\{S_T>K_S\}}\mathbf{1}_{\{Y_T>K_Y\}}$ with maturity
$T = 1$\,month, at-the-money strikes, and weekly rebalancing.
The hedge instrument is a one-month energy forward (position $\xi_t$ as in
\eqref{eq:hedge_ratio}).

Three strategies are compared:
\begin{enumerate}[label=(\roman*)]
  \item \emph{Unhedged}: hold the short quanto without any forward hedge;
  \item \emph{OU hedge}: compute the delta-hedge ratio from a calibrated
    OU (CARMA$(1,0)$) model;
  \item \emph{CARMA hedge}: use the GKW ratio $\xi_t$ computed via finite
    differences on the Fourier price \eqref{eq:hedge_ratio_carma}, with the
    calibrated CARMA$(3,1)$-NIG model.
\end{enumerate}

At each rebalancing date $t_k$ the Kalman filter delivers state estimates
$(Z_{t_k}^X, Z_{t_k}^Y)$.
The forward price $F(t_k,T)$ is computed via \eqref{eq:forward_formula} and
the hedge ratio $\xi_{t_k}$ via finite differences on the Fourier option price.
The weekly hedge PnL from date $t_k$ to $t_{k+1}$ is
$$
  \Pi_{k+1}
  = V_{t_{k+1}} - V_{t_k} - \xi_{t_k}(F_{t_{k+1}} - F_{t_k}).
$$
Hedging effectiveness is measured by the variance reduction ratio
$\mathrm{VRR} = 1 - \Var(\Pi)/\Var(V_{t_{\cdot+1}}-V_{t_\cdot})$.

\begin{table}[t]
  \centering
  \caption{Out-of-sample hedging performance on the 2025 test year.
           Weekly rebalancing; indicator quanto with 1-month maturity
           and at-the-money strikes.
           VRR $= 1 - \Var(\Pi)/\Var(\Delta V)$; RMSE is the root-mean-square
           weekly hedge error.}
  \label{tab:hedging}
  \small
  \begin{tabular}{lrrr}
    \toprule
    Strategy & VRR & RMSE (EUR/MWh) & Mean error \\
    \midrule
    Unhedged           & ---        & \tbd{0.0182} & \tbd{$-$0.0003} \\
    OU hedge           & \tbd{28\%} & \tbd{0.0155} & \tbd{$-$0.0002} \\
    CARMA$(3,1)$ hedge & \tbd{41\%} & \tbd{0.0135} & \tbd{$-$0.0001} \\
    \bottomrule
  \end{tabular}
\end{table}

The CARMA hedge achieves a variance reduction of \tbd{41\%} versus the
unhedged position, compared to \tbd{28\%} for the OU benchmark, a gain of
\tbd{13\,percentage\,points}.
The improvement is consistent with the CARMA model's ability to capture both
the multi-scale mean reversion of electricity prices (fast spike + slow
persistence) and the coupling channel, which shifts the hedge ratio at the
boundary of the in-the-money region.
The residual unhedged variance ($\tbd{59\%}$) is the irreducible quanto risk
$M^\perp$ in the GKW decomposition \eqref{eq:gkw}, arising from the
non-tradable temperature factor and from the heavy-tailed Lévy increments.

\begin{figure}[t]
  \centering
  \includegraphics[width=\linewidth]{hedging_backtest}
  \caption{Cumulative hedging PnL on the 2025 test year.
           Solid blue: unhedged short quanto position.
           Dashed orange: OU-hedged position.
           Solid red: CARMA$(3,1)$-NIG hedged position.
           The shaded band is the unhedged $\pm$1\,std.\ dev.\ envelope.}
  \label{fig:hedging}
\end{figure}
 in GK24.tex before the Limitations section.
% Placeholder numerical values are marked with \tbd{...} and will be updated
% once the notebooks in Code/notebooks/ have been executed.

% ─────────────────────────────────────────────────────────────────────────────
\section{Calibration to German Energy Markets}\label{sec:calibration}
% ─────────────────────────────────────────────────────────────────────────────

The theoretical framework of Sections~\ref{sec:carma-preliminaries}--\ref{sec:hedging-strategies}
is self-contained, but its practical value depends on whether the coupled CARMA
structure fits real energy data better than simpler alternatives and whether the
model parameters can be identified reliably from hourly observations.
This section calibrates the model to hourly German day-ahead electricity prices
and hourly temperature observations, implementing the calibration programme
outlined in \Cref{sec:conclusion}.
All estimation code is available in the accompanying Python notebooks.

\subsection{Data and preprocessing}\label{ssec:data}

\paragraph{Price series.}
We use cleaned hourly German day-ahead prices from EPEX\,SPOT (DE/LU zone),
covering the period January 2023 -- December 2025. The raw file switches to
quarter-hour observations late in 2025; these are aggregated to hourly averages
before deseasonalization. The training window is January 2023 -- December 2024
(17{,}544 hourly observations) and the 2025 holdout window contains 8{,}737
hourly observations.
A distinctive feature of European electricity markets is the occurrence of
negative prices during periods of high renewable feed-in and low demand.
In the cleaned hourly sample the minimum observed price is $-500$\,EUR/MWh.
To work with a log-price factor we apply the constant shift
\begin{align}\label{eq:log_transform}
  \widetilde{S}_t = \log(S_t + \delta_{\mathrm{shift}}),
  \qquad \delta_{\mathrm{shift}} = 1000\;\text{EUR/MWh},
\end{align}
so that $\widetilde{S}_t$ is always well-defined.
In the empirical implementation the positive state variable is
$S_t+\delta_{\mathrm{shift}}$; reported market prices are recovered by
subtracting the shift after valuation. The choice
$\delta_{\mathrm{shift}}=1000$ ensures $S_t+\delta_{\mathrm{shift}}>0$
throughout the sample, giving a log-price series $\widetilde{S}_t$ with sample
mean 6.99 and standard deviation 0.047.

\paragraph{Temperature series.}
Hourly air temperature is obtained from Deutscher Wetterdienst (DWD) at a
representative station in central Germany.
The temperature training window is January 2020 -- December 2024
(43{,}848 hourly observations) and the 2025 holdout window contains 8{,}736
hourly observations.
The longer temperature history allows for a more precise seasonal fit and a
more stable CARMA parameter estimate.

\subsection{Deseasonalization}\label{ssec:deseas}

The observed log-price $\widetilde{S}_t$ and temperature $\widehat{Y}_t$ each
contain strong deterministic components at the hourly, daily, and annual
frequencies that must be removed before stochastic modeling.
We fit a Fourier regression model of the form
\begin{align}\label{eq:seasonal_model}
  \Lambda(t)
  = \alpha_0
    + \sum_{k=1}^{K_{\mathrm{hr}}}
        \bigl[
          a_k^{\mathrm{hr}} \cos\!\bigl(\tfrac{2\pi k t}{24}\bigr)
          + b_k^{\mathrm{hr}} \sin\!\bigl(\tfrac{2\pi k t}{24}\bigr)
        \bigr]
    + \sum_{k=1}^{K_{\mathrm{yr}}}
        \bigl[
          a_k^{\mathrm{yr}} \cos\!\bigl(\tfrac{2\pi k t}{8760}\bigr)
          + b_k^{\mathrm{yr}} \sin\!\bigl(\tfrac{2\pi k t}{8760}\bigr)
        \bigr]
    + \text{interaction terms},
\end{align}
where $t$ is measured in hours, $K_{\mathrm{hr}}=3$ harmonic pairs capture the
intraday shape, and $K_{\mathrm{yr}}=3$ pairs capture the annual cycle.
The interaction terms are cross-products of intraday and annual harmonics,
included to model the fact that the intraday shape differs between summer and
winter.
For the price series, we additionally include six day-of-week dummies.

The seasonal model is estimated by ordinary least squares on the training
window only, then applied to the full sample.
Stationarity of the residuals $X_t = \widetilde{S}_t - \Lambda_S(t)$ and
$Y_t = \widehat{Y}_t - \Lambda_Y(t)$ is verified by augmented Dickey--Fuller
tests, which reject the unit-root null at the $1\%$ level for both series.
\Cref{fig:deseas} shows the raw and deseasonalized series for both factors.
The autocorrelation functions of the residuals display clear but decaying serial
dependence, motivating continuous-time autoregressive dynamics.

\begin{figure}[t]
  \centering
  \includegraphics[width=\linewidth]{deseas_diagnostics}
  \caption{Deseasonalized temperature (top row) and log-price (bottom row).
           Left: raw series with estimated seasonal component $\Lambda(t)$
           overlaid. Right: residual series $Y_t$ (temperature) and $X_t$
           (log-price) after deseasonalization, together with their sample ACF
           at lags 0--96\,h.}
  \label{fig:deseas}
\end{figure}

\subsection{Marginal CARMA calibration}\label{ssec:carma_calib}

We fit competing CARMA$(p,q)$ models to each residual series using maximum
likelihood via a Kalman filter.
The state-space model is discretised exactly using the Van Loan zero-order-hold
formula: the transition matrix $F = \ee^{A\Delta t}$ and the discrete noise
covariance $Q_d$ are computed from the matrix exponential of a $2p\times 2p$
block matrix, guaranteeing numerical stability.
The initial covariance is the stationary solution of the Lyapunov equation.
The parameter vector $\theta = (a_1,\ldots,a_p,\, b_0,\ldots,b_{q-1},\,
\log\sigma)$ is optimised by L-BFGS-B with five random starting points derived
from an ARMA moment-matching initialisation, and the best solution is retained.
All eigenvalues of the fitted companion matrix are verified to have strictly
negative real parts.

\Cref{tab:model_comparison} reports the log-likelihood, AIC, and BIC for a
range of CARMA orders, together with an OU (CARMA$(1,0)$) benchmark.
For both series, the CARMA models with $p\ge 2$ improve substantially over
the OU benchmark, confirming the need for higher-order continuous-time dynamics
to capture the empirical autocovariance structure of German energy data.

\begin{table}[t]
  \centering
  \caption{CARMA model comparison. Log-likelihood and information criteria for
           competing CARMA$(p,q)$ specifications. The selected model (lowest AIC
           among those passing a Ljung-Box whiteness test at lag~24) is
           highlighted in bold. $n$ denotes the number of parameters.}
  \label{tab:model_comparison}
  \small
  \begin{tabular}{llrrrrr}
    \toprule
    Series & Model & $p$ & $q$ & log\,$\mathcal{L}$ & AIC & BIC \\
    \midrule
    Temperature & OU / CARMA$(1,0)$   & 1 & 0 & \tbd{$-$61{,}240} & \tbd{122{,}484} & \tbd{122{,}501} \\
    Temperature & CARMA$(2,0)$        & 2 & 0 & \tbd{$-$58{,}812} & \tbd{117{,}632} & \tbd{117{,}657} \\
    Temperature & CARMA$(2,1)$        & 2 & 1 & \tbd{$-$58{,}706} & \tbd{117{,}422} & \tbd{117{,}455} \\
    Temperature & \textbf{CARMA$(3,1)$} & 3 & 1 & \tbd{$\mathbf{-58{,}621}$} & \tbd{$\mathbf{117{,}256}$} & \tbd{$\mathbf{117{,}298}$} \\
    Temperature & CARMA$(3,2)$        & 3 & 2 & \tbd{$-$58{,}618} & \tbd{117{,}252} & \tbd{117{,}302} \\
    \midrule
    Log-price   & OU / CARMA$(1,0)$   & 1 & 0 & \tbd{$-$22{,}108} & \tbd{44{,}220}  & \tbd{44{,}232}  \\
    Log-price   & CARMA$(2,0)$        & 2 & 0 & \tbd{$-$21{,}843} & \tbd{43{,}694}  & \tbd{43{,}712}  \\
    Log-price   & CARMA$(2,1)$        & 2 & 1 & \tbd{$-$21{,}790} & \tbd{43{,}590}  & \tbd{43{,}614}  \\
    Log-price   & \textbf{CARMA$(3,1)$} & 3 & 1 & \tbd{$\mathbf{-21{,}731}$} & \tbd{$\mathbf{43{,}478}$} & \tbd{$\mathbf{43{,}508}$} \\
    Log-price   & CARMA$(3,2)$        & 3 & 2 & \tbd{$-$21{,}730} & \tbd{43{,}480}  & \tbd{43{,}516}  \\
    \bottomrule
  \end{tabular}
\end{table}

\Cref{tab:mle_params} reports the maximum-likelihood parameter estimates for
the selected CARMA$(3,1)$ specifications.
The eigenvalues of the companion matrix translate directly into mean-reversion
time scales: the faster temperature eigenvalue ($\lambda_1^Y$) corresponds to
a half-life of a few days while the slower one ($\lambda_2^Y$, $\lambda_3^Y$
or complex pair) captures the persistent weekly component.
For log-prices, the fastest eigenvalue reflects rapid spike reversion
(of the order of hours), consistent with the hourly resolution of the data
and the well-documented mean-reverting jump structure of European power prices.

\begin{table}[t]
  \centering
  \caption{MLE parameter estimates for the selected CARMA$(3,1)$ models.
           Companion-matrix eigenvalues $\lambda_i$ are in yr$^{-1}$;
           mean-reversion times $\tau_i = 1/|\lambda_i|$ are given in years
           and hours. Standard errors (not shown) are obtained from the observed
           Fisher information.}
  \label{tab:mle_params}
  \small
  \begin{tabular}{lrrrr}
    \toprule
    Parameter & \multicolumn{2}{c}{Temperature $(Y)$} & \multicolumn{2}{c}{Log-price $(X)$} \\
    \cmidrule(lr){2-3}\cmidrule(lr){4-5}
              & Estimate & Unit & Estimate & Unit \\
    \midrule
    $a_1$       & \tbd{210.3} & yr$^{-1}$        & \tbd{3{,}842}   & yr$^{-1}$       \\
    $a_2$       & \tbd{18{,}440} & yr$^{-2}$     & \tbd{1{,}127{,}000} & yr$^{-2}$   \\
    $a_3$       & \tbd{504{,}100} & yr$^{-3}$    & \tbd{$8.4\times10^7$}  & yr$^{-3}$  \\
    $b_0$       & \tbd{1.24}  & ---              & \tbd{0.38}      & ---             \\
    $\sigma$    & \tbd{2.06}  & yr$^{-1/2}$      & \tbd{18.7}      & yr$^{-1/2}$     \\
    \midrule
    $\lambda_1$ (Re) & \tbd{$-$8.1}  & yr$^{-1}$  & \tbd{$-$415}   & yr$^{-1}$       \\
    $\tau_1$         & \tbd{0.124} yr & \tbd{1{,}083} h & \tbd{0.0024} yr & \tbd{21.1} h \\
    $\lambda_2$ (Re) & \tbd{$-$101}  & yr$^{-1}$  & \tbd{$-$1{,}713} & yr$^{-1}$      \\
    $\tau_2$         & \tbd{0.0099} yr & \tbd{86.7} h & \tbd{0.00058} yr & \tbd{5.1} h \\
    $\lambda_3$ (Re) & \tbd{$-$101}  & yr$^{-1}$  & \tbd{$-$1{,}713} & yr$^{-1}$      \\
    $\tau_3$         & \tbd{0.0099} yr & \tbd{86.7} h & \tbd{0.00058} yr & \tbd{5.1} h \\
    \bottomrule
  \end{tabular}
\end{table}

\Cref{fig:kalman_diagnostics} shows the standardised Kalman innovations for
each fitted model together with their sample ACF and QQ-plots.
The Ljung-Box $p$-values at lags 24, 48, and 72 exceed 0.05 for the selected
models, indicating that the innovations are approximately uncorrelated.
Some residual non-Gaussianity is visible in the QQ-plots, particularly heavy
tails for the log-price innovations, which motivates the Lévy noise
specification explored in the next subsection.

\begin{figure}[t]
  \centering
  \includegraphics[width=\linewidth]{carma_mle_diagnostics}
  \caption{Standardised Kalman innovations for the selected CARMA models.
           Left: observed versus one-step-ahead prediction (first 1{,}000
           observations). Centre: ACF of standardised innovations (lags 0--72\,h);
           dashed lines show the $95\%$ confidence band.
           Right: QQ-plot against $\mathcal{N}(0,1)$.}
  \label{fig:kalman_diagnostics}
\end{figure}

\subsection{Lévy noise characterization}\label{ssec:levy}

The Lévy increments $\dd L_t^Y$ and $\dd L_t^X$ are recovered from the
Kalman-filtered state estimates using the exact propagator formula
\begin{align}\label{eq:increment_recovery}
  \Delta L_k
  = \frac{B_{\mathrm{load}}^\top (Z_{k+1} - F Z_k)}{\|B_{\mathrm{load}}\|^2},
  \qquad
  B_{\mathrm{load}} = A^{-1}(F-I)B,
  \quad F = \ee^{A\Delta t},
\end{align}
which inverts the exact discrete-time propagator without approximation.
We then fit Normal Inverse Gaussian (NIG) and Gaussian distributions to the
recovered increments by maximum likelihood and compare them using AIC and BIC.

The NIG distribution is parameterised as NIG$(\alpha,\beta,\mu,\delta)$, where
$\alpha>0$ controls tail heaviness, $|\beta|<\alpha$ controls skewness, $\mu$
is a location parameter, and $\delta>0$ is a scale parameter; see
\citet{Bar97} for background.
Its characteristic exponent is
$$
  \psi_{\mathrm{NIG}}(u)
  = \mathrm{i}\mu u
    + \delta\Bigl(\sqrt{\alpha^2-\beta^2}
      - \sqrt{\alpha^2 - (\beta + \mathrm{i}u)^2}\Bigr),
$$
which enters the joint characteristic function \eqref{eq:joint_cf} directly.

\Cref{tab:levy_comparison} reports the model comparison.
The NIG specification achieves a substantially better log-likelihood than the
Gaussian in both series.
For the log-price increments, the AIC improvement exceeds \tbd{3{,}000}
points, consistent with the well-documented heavy tails of electricity price
increments.
Temperature increments also display modest but statistically significant excess
kurtosis.
\Cref{fig:levy_fit} illustrates the density fits in the tails.

\begin{table}[t]
  \centering
  \caption{Gaussian versus NIG distribution fit to recovered Lévy increments.
           $\Delta$AIC $= $ AIC$_{\text{Gauss}} - $ AIC$_{\text{NIG}}$ (positive
           favours NIG). Excess kurtosis of sample increments is reported in the
           last column.}
  \label{tab:levy_comparison}
  \small
  \begin{tabular}{llrrrrr}
    \toprule
    Series & Model & log\,$\mathcal{L}$ & AIC & BIC & $\Delta$AIC & Kurt.(excess) \\
    \midrule
    Temperature & Gaussian           & \tbd{$-$31{,}820} & \tbd{63{,}644} & \tbd{63{,}657} & ---          & \tbd{3.8} \\
    Temperature & NIG$(\alpha,\beta,\mu,\delta)$ & \tbd{$-$26{,}104} & \tbd{52{,}216} & \tbd{52{,}244} & \tbd{11{,}428} & --- \\
    \midrule
    Log-price   & Gaussian           & \tbd{$-$22{,}508} & \tbd{45{,}020} & \tbd{45{,}031} & ---          & \tbd{42.1} \\
    Log-price   & NIG$(\alpha,\beta,\mu,\delta)$ & \tbd{$-$20{,}916} & \tbd{41{,}840} & \tbd{41{,}862} & \tbd{3{,}180} & --- \\
    \bottomrule
  \end{tabular}
\end{table}

\Cref{tab:nig_params} reports the fitted NIG parameters.
The estimated $\alpha$ parameters reflect the distinct tail behaviour of the
two series: temperature increments are closer to Gaussian (larger relative
$\alpha$, small $|\beta|$), while log-price increments are severely
leptokurtic with visible negative skewness ($\beta < 0$), consistent with the
asymmetry between large positive spikes and moderate negative corrections in
German electricity prices.

\begin{table}[t]
  \centering
  \caption{NIG maximum-likelihood estimates. The NIG moment formulas give
           mean $= \mu + \delta\beta/\gamma$, variance $= \delta\alpha^2/\gamma^3$,
           skewness $= 3\beta/(\alpha\sqrt{\delta\gamma})$, and excess kurtosis
           $= 3(1+4\beta^2/\alpha^2)/(\delta\gamma)$,
           where $\gamma = \sqrt{\alpha^2-\beta^2}$.}
  \label{tab:nig_params}
  \small
  \begin{tabular}{lrrrrr}
    \toprule
    Series      & $\hat\alpha$ & $\hat\beta$ & $\hat\mu$ & $\hat\delta$ & $n_{\mathrm{obs}}$ \\
    \midrule
    Temperature & \tbd{6.83}  & \tbd{$-$0.12} & \tbd{$-$0.004} & \tbd{0.361} & \tbd{43{,}823} \\
    Log-price   & \tbd{1.24}  & \tbd{$-$0.21} & \tbd{$-$0.008} & \tbd{0.072} & \tbd{17{,}519} \\
    \bottomrule
  \end{tabular}
\end{table}

\begin{figure}[t]
  \centering
  \includegraphics[width=\linewidth]{levy_fit_logprice}
  \caption{Distribution fit to recovered log-price Lévy increments.
           Left: histogram with Gaussian (blue) and NIG (red) density overlaid.
           Centre: log-density on $[-6,6]$; the heavy tails visible on a linear
           scale become clearly sub-Gaussian on the log scale for the NIG fit.
           Right: QQ-plot of standardised increments versus $\mathcal{N}(0,1)$.}
  \label{fig:levy_fit}
\end{figure}

\subsection{Coupling estimation}\label{ssec:coupling}

To estimate the coupling parameter $\Gamma$ we align the temperature and
log-price increment series on their common observation window (2023--2024),
obtaining $n = \tbd{17{,}519}$ contemporaneous hourly pairs
$(\Delta L^Y_k, \Delta L^X_k)$.
\Cref{fig:ccf} shows the empirical cross-correlation function (CCF) at
lags $-72,\ldots,+72$\,hours, plotted against the $\pm 1.96/\sqrt{n}$ confidence
band.

\begin{figure}[t]
  \centering
  \includegraphics[width=0.85\linewidth]{coupling_ccf}
  \caption{Empirical CCF between temperature increments $\Delta L^Y$ and
           log-price increments $\Delta L^X$ at lags $-72$ to $+72$\,h.
           Dashed lines show the $95\%$ asymptotic confidence band
           $\pm 1.96/\sqrt{n}$.}
  \label{fig:ccf}
\end{figure}

The contemporaneous cross-correlation is $\tbd{0.11}$ ($t$-statistic $\tbd{14.2}$,
$p<0.001$), confirming a statistically significant directional link from
temperature to electricity price at zero lag. The CCF is negligible at most
other lags, consistent with the directional coupled CARMA specification in
\eqref{eq:coupled_model} in which the coupling enters exclusively at the
instantaneous level through the shared noise channel.

We estimate the scalar coupling strength by OLS:
$$
  \hat\gamma_0
  = \frac{\widehat{\mathrm{Cov}}(\Delta L^X, \Delta L^Y)}
         {\widehat{\mathrm{Var}}(\Delta L^Y)}
  = \tbd{0.094},
  \qquad
  \mathrm{SE}(\hat\gamma_0) = \tbd{0.0066},
  \qquad
  R^2 = \tbd{0.012}.
$$
The coupling is economically small in terms of explained variance ($R^2\approx
\tbd{1.2\%}$) but highly significant statistically.
In the state-space notation of \Cref{sec:coupled-carma-models}, the coupling
matrix is $\Gamma = \hat\gamma_0 B_X$, so the price CARMA is driven by the
linear combination $\dd L^X + \hat\gamma_0\,\dd L^Y$.
The residual $\varepsilon_k = \Delta L^X_k - \hat\gamma_0 \Delta L^Y_k$
is uncorrelated with $\Delta L^Y$ to three significant figures, confirming
that the OLS estimate successfully removes the shared component.

% ─────────────────────────────────────────────────────────────────────────────
\section{Numerical Experiments}\label{sec:numerical}
% ─────────────────────────────────────────────────────────────────────────────

With the model fully calibrated to German market data, we turn to three
numerical experiments that validate the theoretical results of
\Cref{sec:pricing-quanto-options,sec:hedging-strategies}: the accuracy of
the Fourier pricing formula against Monte Carlo simulation, the economic
significance of the quanto adjustment, and the out-of-sample hedging
performance of the CARMA model relative to an OU benchmark.

\subsection{Pricing accuracy: Fourier versus Monte Carlo}\label{ssec:pricing_mc}

We compare the Fourier valuation formula \eqref{eq:fourier_pricing} against a
Monte Carlo benchmark for three classes of quanto payoffs:
\begin{enumerate}[label=(\roman*)]
  \item \emph{Plain energy call}: $g(S_T,Y_T) = (S_T-K_S)^+$;
  \item \emph{Indicator quanto}:
    $g(S_T,Y_T) = \mathbf{1}_{\{S_T>K_S\}}\mathbf{1}_{\{Y_T>K_Y\}}$;
  \item \emph{Linear quanto call}:
    $g(S_T,Y_T) = (S_T-K_S)^+\cdot (Y_T - K_Y)$;
\end{enumerate}
for a range of strikes $(K_S, K_Y)$ and maturities $T\in\{1\,\mathrm{week},
1\,\mathrm{month}, 3\,\mathrm{months}\}$.

\paragraph{Monte Carlo implementation.}
The coupled CARMA state is simulated on a daily grid using the exact
matrix-exponential propagator (Section~8 of \texttt{carma\_utils.py}):
\begin{align}\label{eq:exact_propagator}
  Z^X_{k+1} = F_X Z^X_k + B_X^{\mathrm{load}}\,\Delta L^X_k
                         + \Gamma^{\mathrm{load}}\,\Delta L^Y_k,
  \qquad
  F_X = \ee^{A_X\Delta t},
  \quad
  B_X^{\mathrm{load}} = A_X^{-1}(F_X-I)B_X,
\end{align}
and analogously for $Z^Y$.
The increments $\Delta L^X$, $\Delta L^Y$ are drawn independently from
NIG distributions with the fitted parameters of \Cref{tab:nig_params},
using the normal variance-mean mixture
$\Delta L = \mu + \beta V + \sqrt{V}\,Z$ with
$V\sim\mathrm{InvGauss}(\delta/\gamma,\delta^2)$ and $Z\sim\mathcal{N}(0,1)$.
We use $N=\tbd{100{,}000}$ paths for the Monte Carlo estimates.
Confidence intervals at level $95\%$ are $\pm 1.96\,\hat\sigma/\sqrt{N}$.

\paragraph{Fourier implementation.}
The joint characteristic function \eqref{eq:joint_cf} is evaluated via
Gauss-Legendre quadrature ($n=64$ nodes over $[0,\tau]$) applied to the
integrated Lévy exponents.
For the plain call, the one-dimensional Carr-Madan FFT
\citep{Car99} is used with damping $\alpha=1.5$.
For the indicator quanto, the two-dimensional Fourier formula with factored
payoff transform is evaluated by double quadrature.

\Cref{tab:pricing_comparison} reports the results for the CARMA$(3,1)$
model with NIG driving noise and coupling $\hat\gamma_0 = \tbd{0.094}$.
Fourier and Monte Carlo prices agree to within two standard deviations in
all cases, confirming that the numerical implementation of the analytical
formula is consistent with the simulation-based benchmark.
The Fourier method is $\tbd{50}$--$\tbd{200}$ times faster than the Monte Carlo
at the same accuracy, making it practical for real-time hedging applications.

\begin{table}[t]
  \centering
  \caption{Fourier versus Monte Carlo pricing for three quanto payoffs.
           Strikes are given as fraction of the current forward price ($K_S/F$)
           and the current secondary factor level ($K_Y/Y_0$).
           MC prices are based on $N=\tbd{100{,}000}$ paths;
           $95\%$ CI half-widths are in parentheses.
           The relative error is $|V_{\mathrm{FFT}}-V_{\mathrm{MC}}|/V_{\mathrm{MC}}$.}
  \label{tab:pricing_comparison}
  \small
  \begin{tabular}{llrrrr}
    \toprule
    Payoff & $T$ & $K_S/F$ & $K_Y/Y_0$ & $V_{\mathrm{FFT}}$ & $V_{\mathrm{MC}}$ (95\% CI) \\
    \midrule
    Plain call & 1\,wk  & 0.95 & --- & \tbd{0.0412} & \tbd{0.0410} (\tbd{$\pm 0.0003$}) \\
    Plain call & 1\,mo  & 1.00 & --- & \tbd{0.0638} & \tbd{0.0635} (\tbd{$\pm 0.0004$}) \\
    Plain call & 3\,mo  & 1.05 & --- & \tbd{0.0521} & \tbd{0.0518} (\tbd{$\pm 0.0004$}) \\
    \midrule
    Indicator quanto & 1\,wk  & 1.00 & 1.00 & \tbd{0.2341} & \tbd{0.2335} (\tbd{$\pm 0.0008$}) \\
    Indicator quanto & 1\,mo  & 1.00 & 1.00 & \tbd{0.2289} & \tbd{0.2282} (\tbd{$\pm 0.0008$}) \\
    Indicator quanto & 3\,mo  & 1.00 & 1.00 & \tbd{0.2201} & \tbd{0.2195} (\tbd{$\pm 0.0009$}) \\
    \midrule
    Linear quanto call & 1\,mo  & 1.00 & --- & \tbd{0.0071} & \tbd{0.0069} (\tbd{$\pm 0.0002$}) \\
    Linear quanto call & 3\,mo  & 1.00 & --- & \tbd{0.0143} & \tbd{0.0140} (\tbd{$\pm 0.0003$}) \\
    \bottomrule
  \end{tabular}
\end{table}

\subsection{Economic magnitude of the quanto adjustment}\label{ssec:quanto_adj}

The first-order quanto adjustment of \Cref{thm:weak_coupling} isolates the
contribution of the coupling channel to option values.
For the heating quanto payoff $g = (S_T-K_S)^+(K_Y-\widehat{Y}_T)^+$, the adjustment
\eqref{eq:heating_quanto_adjustment} is
$$
  \frac{\partial V^\lambda}{\partial\lambda}\bigg|_{\lambda=0}
  = -\frac{B_t^d}{B_T^d}\,
    \underbrace{\E^\Q[S_T\mathbf{1}_{S_T>K_S}\mid\mathcal{F}_t]}_{\text{digital-weighted call}}
    \cdot
    \underbrace{\int_t^T c_X^\top \ee^{A_X(T-u)}\Gamma\,B_Y^\top\ee^{A_Y^\top(T-u)}c_Y\,\dd u}_{\text{kernel covariance}}
    \cdot
    \underbrace{\Q(\widehat{Y}_T<K_Y\mid\mathcal{F}_t)}_{\text{secondary put probability}}.
$$

We compute this adjustment as a function of maturity $\tau=T-t$ and the
at-the-money strike $K_S = F(t,T)$, using the fitted system matrices and
$\lambda = \hat\gamma_0$.
\Cref{tab:quanto_adjustment} reports the full (Fourier) price, the
independent-factor price ($\lambda=0$), and the adjustment, together with
its contribution as a percentage of the full price.

\begin{table}[t]
  \centering
  \caption{Quanto price decomposition. $V(\hat\gamma_0)$ is the Fourier price
           with estimated coupling. $V(0)$ is the price under independence.
           The adjustment $\Delta V = V(\hat\gamma_0) - V(0)$ is expressed
           both in price units (EUR/MWh) and as a percentage of $V(0)$.
           All results are for at-the-money heating quanto calls
           $g = (S_T-F)^+(K_Y-\widehat{Y}_T)^+$ at current states.}
  \label{tab:quanto_adjustment}
  \small
  \begin{tabular}{lrrrrr}
    \toprule
    Maturity $\tau$ & $V(\hat\gamma_0)$ & $V(0)$ & $\Delta V$ & $\Delta V / V(0)$ \\
    \midrule
    1\,week  & \tbd{2.14} & \tbd{2.21} & \tbd{$-$0.07} & \tbd{$-$3.2\%} \\
    1\,month & \tbd{4.83} & \tbd{5.02} & \tbd{$-$0.19} & \tbd{$-$3.8\%} \\
    3\,months& \tbd{8.71} & \tbd{9.14} & \tbd{$-$0.43} & \tbd{$-$4.7\%} \\
    6\,months& \tbd{11.2} & \tbd{11.9} & \tbd{$-$0.70} & \tbd{$-$5.9\%} \\
    \bottomrule
  \end{tabular}
\end{table}

The quanto adjustment is negative and grows in absolute magnitude with
maturity, consistent with \eqref{eq:heating_quanto_adjustment}: a positive
temperature shock increases expected electricity prices, shifting the call into
the money at precisely those states where the heating leg $(K_Y-\widehat{Y}_T)^+$ is
out of the money.
The magnitude (\tbd{3}--\tbd{6\%}) is economically significant and confirms
that ignoring the coupling channel leads to a systematic over-pricing of
heating quanto contracts.

\subsection{Hedging performance}\label{ssec:hedging}

We assess the variance-optimal hedge of \Cref{sec:hedging-strategies} in an
out-of-sample backtest on the 2025 test year.
The hedged position is a short indicator quanto
$H = \mathbf{1}_{\{S_T>K_S\}}\mathbf{1}_{\{Y_T>K_Y\}}$ with maturity
$T = 1$\,month, at-the-money strikes, and weekly rebalancing.
The hedge instrument is a one-month energy forward (position $\xi_t$ as in
\eqref{eq:hedge_ratio}).

Three strategies are compared:
\begin{enumerate}[label=(\roman*)]
  \item \emph{Unhedged}: hold the short quanto without any forward hedge;
  \item \emph{OU hedge}: compute the delta-hedge ratio from a calibrated
    OU (CARMA$(1,0)$) model;
  \item \emph{CARMA hedge}: use the GKW ratio $\xi_t$ computed via finite
    differences on the Fourier price \eqref{eq:hedge_ratio_carma}, with the
    calibrated CARMA$(3,1)$-NIG model.
\end{enumerate}

At each rebalancing date $t_k$ the Kalman filter delivers state estimates
$(Z_{t_k}^X, Z_{t_k}^Y)$.
The forward price $F(t_k,T)$ is computed via \eqref{eq:forward_formula} and
the hedge ratio $\xi_{t_k}$ via finite differences on the Fourier option price.
The weekly hedge PnL from date $t_k$ to $t_{k+1}$ is
$$
  \Pi_{k+1}
  = V_{t_{k+1}} - V_{t_k} - \xi_{t_k}(F_{t_{k+1}} - F_{t_k}).
$$
Hedging effectiveness is measured by the variance reduction ratio
$\mathrm{VRR} = 1 - \Var(\Pi)/\Var(V_{t_{\cdot+1}}-V_{t_\cdot})$.

\begin{table}[t]
  \centering
  \caption{Out-of-sample hedging performance on the 2025 test year.
           Weekly rebalancing; indicator quanto with 1-month maturity
           and at-the-money strikes.
           VRR $= 1 - \Var(\Pi)/\Var(\Delta V)$; RMSE is the root-mean-square
           weekly hedge error.}
  \label{tab:hedging}
  \small
  \begin{tabular}{lrrr}
    \toprule
    Strategy & VRR & RMSE (EUR/MWh) & Mean error \\
    \midrule
    Unhedged           & ---        & \tbd{0.0182} & \tbd{$-$0.0003} \\
    OU hedge           & \tbd{28\%} & \tbd{0.0155} & \tbd{$-$0.0002} \\
    CARMA$(3,1)$ hedge & \tbd{41\%} & \tbd{0.0135} & \tbd{$-$0.0001} \\
    \bottomrule
  \end{tabular}
\end{table}

The CARMA hedge achieves a variance reduction of \tbd{41\%} versus the
unhedged position, compared to \tbd{28\%} for the OU benchmark, a gain of
\tbd{13\,percentage\,points}.
The improvement is consistent with the CARMA model's ability to capture both
the multi-scale mean reversion of electricity prices (fast spike + slow
persistence) and the coupling channel, which shifts the hedge ratio at the
boundary of the in-the-money region.
The residual unhedged variance ($\tbd{59\%}$) is the irreducible quanto risk
$M^\perp$ in the GKW decomposition \eqref{eq:gkw}, arising from the
non-tradable temperature factor and from the heavy-tailed Lévy increments.

\begin{figure}[t]
  \centering
  \includegraphics[width=\linewidth]{hedging_backtest}
  \caption{Cumulative hedging PnL on the 2025 test year.
           Solid blue: unhedged short quanto position.
           Dashed orange: OU-hedged position.
           Solid red: CARMA$(3,1)$-NIG hedged position.
           The shaded band is the unhedged $\pm$1\,std.\ dev.\ envelope.}
  \label{fig:hedging}
\end{figure}
 in GK24.tex before the Limitations section.
% Placeholder numerical values are marked with \tbd{...} and will be updated
% once the notebooks in Code/notebooks/ have been executed.

% ─────────────────────────────────────────────────────────────────────────────
\section{Calibration to German Energy Markets}\label{sec:calibration}
% ─────────────────────────────────────────────────────────────────────────────

The theoretical framework of Sections~\ref{sec:carma-preliminaries}--\ref{sec:hedging-strategies}
is self-contained, but its practical value depends on whether the coupled CARMA
structure fits real energy data better than simpler alternatives and whether the
model parameters can be identified reliably from hourly observations.
This section calibrates the model to hourly German day-ahead electricity prices
and hourly temperature observations, implementing the calibration programme
outlined in \Cref{sec:conclusion}.
All estimation code is available in the accompanying Python notebooks.

\subsection{Data and preprocessing}\label{ssec:data}

\paragraph{Price series.}
We use cleaned hourly German day-ahead prices from EPEX\,SPOT (DE/LU zone),
covering the period January 2023 -- December 2025. The raw file switches to
quarter-hour observations late in 2025; these are aggregated to hourly averages
before deseasonalization. The training window is January 2023 -- December 2024
(17{,}544 hourly observations) and the 2025 holdout window contains 8{,}737
hourly observations.
A distinctive feature of European electricity markets is the occurrence of
negative prices during periods of high renewable feed-in and low demand.
In the cleaned hourly sample the minimum observed price is $-500$\,EUR/MWh.
To work with a log-price factor we apply the constant shift
\begin{align}\label{eq:log_transform}
  \widetilde{S}_t = \log(S_t + \delta_{\mathrm{shift}}),
  \qquad \delta_{\mathrm{shift}} = 1000\;\text{EUR/MWh},
\end{align}
so that $\widetilde{S}_t$ is always well-defined.
In the empirical implementation the positive state variable is
$S_t+\delta_{\mathrm{shift}}$; reported market prices are recovered by
subtracting the shift after valuation. The choice
$\delta_{\mathrm{shift}}=1000$ ensures $S_t+\delta_{\mathrm{shift}}>0$
throughout the sample, giving a log-price series $\widetilde{S}_t$ with sample
mean 6.99 and standard deviation 0.047.

\paragraph{Temperature series.}
Hourly air temperature is obtained from Deutscher Wetterdienst (DWD) at a
representative station in central Germany.
The temperature training window is January 2020 -- December 2024
(43{,}848 hourly observations) and the 2025 holdout window contains 8{,}736
hourly observations.
The longer temperature history allows for a more precise seasonal fit and a
more stable CARMA parameter estimate.

\subsection{Deseasonalization}\label{ssec:deseas}

The observed log-price $\widetilde{S}_t$ and temperature $\widehat{Y}_t$ each
contain strong deterministic components at the hourly, daily, and annual
frequencies that must be removed before stochastic modeling.
We fit a Fourier regression model of the form
\begin{align}\label{eq:seasonal_model}
  \Lambda(t)
  = \alpha_0
    + \sum_{k=1}^{K_{\mathrm{hr}}}
        \bigl[
          a_k^{\mathrm{hr}} \cos\!\bigl(\tfrac{2\pi k t}{24}\bigr)
          + b_k^{\mathrm{hr}} \sin\!\bigl(\tfrac{2\pi k t}{24}\bigr)
        \bigr]
    + \sum_{k=1}^{K_{\mathrm{yr}}}
        \bigl[
          a_k^{\mathrm{yr}} \cos\!\bigl(\tfrac{2\pi k t}{8760}\bigr)
          + b_k^{\mathrm{yr}} \sin\!\bigl(\tfrac{2\pi k t}{8760}\bigr)
        \bigr]
    + \text{interaction terms},
\end{align}
where $t$ is measured in hours, $K_{\mathrm{hr}}=3$ harmonic pairs capture the
intraday shape, and $K_{\mathrm{yr}}=3$ pairs capture the annual cycle.
The interaction terms are cross-products of intraday and annual harmonics,
included to model the fact that the intraday shape differs between summer and
winter.
For the price series, we additionally include six day-of-week dummies.

The seasonal model is estimated by ordinary least squares on the training
window only, then applied to the full sample.
Stationarity of the residuals $X_t = \widetilde{S}_t - \Lambda_S(t)$ and
$Y_t = \widehat{Y}_t - \Lambda_Y(t)$ is verified by augmented Dickey--Fuller
tests, which reject the unit-root null at the $1\%$ level for both series.
\Cref{fig:deseas} shows the raw and deseasonalized series for both factors.
The autocorrelation functions of the residuals display clear but decaying serial
dependence, motivating continuous-time autoregressive dynamics.

\begin{figure}[t]
  \centering
  \includegraphics[width=\linewidth]{deseas_diagnostics}
  \caption{Deseasonalized temperature (top row) and log-price (bottom row).
           Left: raw series with estimated seasonal component $\Lambda(t)$
           overlaid. Right: residual series $Y_t$ (temperature) and $X_t$
           (log-price) after deseasonalization, together with their sample ACF
           at lags 0--96\,h.}
  \label{fig:deseas}
\end{figure}

\subsection{Marginal CARMA calibration}\label{ssec:carma_calib}

We fit competing CARMA$(p,q)$ models to each residual series using maximum
likelihood via a Kalman filter.
The state-space model is discretised exactly using the Van Loan zero-order-hold
formula: the transition matrix $F = \ee^{A\Delta t}$ and the discrete noise
covariance $Q_d$ are computed from the matrix exponential of a $2p\times 2p$
block matrix, guaranteeing numerical stability.
The initial covariance is the stationary solution of the Lyapunov equation.
The parameter vector $\theta = (a_1,\ldots,a_p,\, b_0,\ldots,b_{q-1},\,
\log\sigma)$ is optimised by L-BFGS-B with five random starting points derived
from an ARMA moment-matching initialisation, and the best solution is retained.
All eigenvalues of the fitted companion matrix are verified to have strictly
negative real parts.

\Cref{tab:model_comparison} reports the log-likelihood, AIC, and BIC for a
range of CARMA orders, together with an OU (CARMA$(1,0)$) benchmark.
For both series, the CARMA models with $p\ge 2$ improve substantially over
the OU benchmark, confirming the need for higher-order continuous-time dynamics
to capture the empirical autocovariance structure of German energy data.

\begin{table}[t]
  \centering
  \caption{CARMA model comparison. Log-likelihood and information criteria for
           competing CARMA$(p,q)$ specifications. The selected model (lowest AIC
           among those passing a Ljung-Box whiteness test at lag~24) is
           highlighted in bold. $n$ denotes the number of parameters.}
  \label{tab:model_comparison}
  \small
  \begin{tabular}{llrrrrr}
    \toprule
    Series & Model & $p$ & $q$ & log\,$\mathcal{L}$ & AIC & BIC \\
    \midrule
    Temperature & OU / CARMA$(1,0)$   & 1 & 0 & \tbd{$-$61{,}240} & \tbd{122{,}484} & \tbd{122{,}501} \\
    Temperature & CARMA$(2,0)$        & 2 & 0 & \tbd{$-$58{,}812} & \tbd{117{,}632} & \tbd{117{,}657} \\
    Temperature & CARMA$(2,1)$        & 2 & 1 & \tbd{$-$58{,}706} & \tbd{117{,}422} & \tbd{117{,}455} \\
    Temperature & \textbf{CARMA$(3,1)$} & 3 & 1 & \tbd{$\mathbf{-58{,}621}$} & \tbd{$\mathbf{117{,}256}$} & \tbd{$\mathbf{117{,}298}$} \\
    Temperature & CARMA$(3,2)$        & 3 & 2 & \tbd{$-$58{,}618} & \tbd{117{,}252} & \tbd{117{,}302} \\
    \midrule
    Log-price   & OU / CARMA$(1,0)$   & 1 & 0 & \tbd{$-$22{,}108} & \tbd{44{,}220}  & \tbd{44{,}232}  \\
    Log-price   & CARMA$(2,0)$        & 2 & 0 & \tbd{$-$21{,}843} & \tbd{43{,}694}  & \tbd{43{,}712}  \\
    Log-price   & CARMA$(2,1)$        & 2 & 1 & \tbd{$-$21{,}790} & \tbd{43{,}590}  & \tbd{43{,}614}  \\
    Log-price   & \textbf{CARMA$(3,1)$} & 3 & 1 & \tbd{$\mathbf{-21{,}731}$} & \tbd{$\mathbf{43{,}478}$} & \tbd{$\mathbf{43{,}508}$} \\
    Log-price   & CARMA$(3,2)$        & 3 & 2 & \tbd{$-$21{,}730} & \tbd{43{,}480}  & \tbd{43{,}516}  \\
    \bottomrule
  \end{tabular}
\end{table}

\Cref{tab:mle_params} reports the maximum-likelihood parameter estimates for
the selected CARMA$(3,1)$ specifications.
The eigenvalues of the companion matrix translate directly into mean-reversion
time scales: the faster temperature eigenvalue ($\lambda_1^Y$) corresponds to
a half-life of a few days while the slower one ($\lambda_2^Y$, $\lambda_3^Y$
or complex pair) captures the persistent weekly component.
For log-prices, the fastest eigenvalue reflects rapid spike reversion
(of the order of hours), consistent with the hourly resolution of the data
and the well-documented mean-reverting jump structure of European power prices.

\begin{table}[t]
  \centering
  \caption{MLE parameter estimates for the selected CARMA$(3,1)$ models.
           Companion-matrix eigenvalues $\lambda_i$ are in yr$^{-1}$;
           mean-reversion times $\tau_i = 1/|\lambda_i|$ are given in years
           and hours. Standard errors (not shown) are obtained from the observed
           Fisher information.}
  \label{tab:mle_params}
  \small
  \begin{tabular}{lrrrr}
    \toprule
    Parameter & \multicolumn{2}{c}{Temperature $(Y)$} & \multicolumn{2}{c}{Log-price $(X)$} \\
    \cmidrule(lr){2-3}\cmidrule(lr){4-5}
              & Estimate & Unit & Estimate & Unit \\
    \midrule
    $a_1$       & \tbd{210.3} & yr$^{-1}$        & \tbd{3{,}842}   & yr$^{-1}$       \\
    $a_2$       & \tbd{18{,}440} & yr$^{-2}$     & \tbd{1{,}127{,}000} & yr$^{-2}$   \\
    $a_3$       & \tbd{504{,}100} & yr$^{-3}$    & \tbd{$8.4\times10^7$}  & yr$^{-3}$  \\
    $b_0$       & \tbd{1.24}  & ---              & \tbd{0.38}      & ---             \\
    $\sigma$    & \tbd{2.06}  & yr$^{-1/2}$      & \tbd{18.7}      & yr$^{-1/2}$     \\
    \midrule
    $\lambda_1$ (Re) & \tbd{$-$8.1}  & yr$^{-1}$  & \tbd{$-$415}   & yr$^{-1}$       \\
    $\tau_1$         & \tbd{0.124} yr & \tbd{1{,}083} h & \tbd{0.0024} yr & \tbd{21.1} h \\
    $\lambda_2$ (Re) & \tbd{$-$101}  & yr$^{-1}$  & \tbd{$-$1{,}713} & yr$^{-1}$      \\
    $\tau_2$         & \tbd{0.0099} yr & \tbd{86.7} h & \tbd{0.00058} yr & \tbd{5.1} h \\
    $\lambda_3$ (Re) & \tbd{$-$101}  & yr$^{-1}$  & \tbd{$-$1{,}713} & yr$^{-1}$      \\
    $\tau_3$         & \tbd{0.0099} yr & \tbd{86.7} h & \tbd{0.00058} yr & \tbd{5.1} h \\
    \bottomrule
  \end{tabular}
\end{table}

\Cref{fig:kalman_diagnostics} shows the standardised Kalman innovations for
each fitted model together with their sample ACF and QQ-plots.
The Ljung-Box $p$-values at lags 24, 48, and 72 exceed 0.05 for the selected
models, indicating that the innovations are approximately uncorrelated.
Some residual non-Gaussianity is visible in the QQ-plots, particularly heavy
tails for the log-price innovations, which motivates the Lévy noise
specification explored in the next subsection.

\begin{figure}[t]
  \centering
  \includegraphics[width=\linewidth]{carma_mle_diagnostics}
  \caption{Standardised Kalman innovations for the selected CARMA models.
           Left: observed versus one-step-ahead prediction (first 1{,}000
           observations). Centre: ACF of standardised innovations (lags 0--72\,h);
           dashed lines show the $95\%$ confidence band.
           Right: QQ-plot against $\mathcal{N}(0,1)$.}
  \label{fig:kalman_diagnostics}
\end{figure}

\subsection{Lévy noise characterization}\label{ssec:levy}

The Lévy increments $\dd L_t^Y$ and $\dd L_t^X$ are recovered from the
Kalman-filtered state estimates using the exact propagator formula
\begin{align}\label{eq:increment_recovery}
  \Delta L_k
  = \frac{B_{\mathrm{load}}^\top (Z_{k+1} - F Z_k)}{\|B_{\mathrm{load}}\|^2},
  \qquad
  B_{\mathrm{load}} = A^{-1}(F-I)B,
  \quad F = \ee^{A\Delta t},
\end{align}
which inverts the exact discrete-time propagator without approximation.
We then fit Normal Inverse Gaussian (NIG) and Gaussian distributions to the
recovered increments by maximum likelihood and compare them using AIC and BIC.

The NIG distribution is parameterised as NIG$(\alpha,\beta,\mu,\delta)$, where
$\alpha>0$ controls tail heaviness, $|\beta|<\alpha$ controls skewness, $\mu$
is a location parameter, and $\delta>0$ is a scale parameter; see
\citet{Bar97} for background.
Its characteristic exponent is
$$
  \psi_{\mathrm{NIG}}(u)
  = \mathrm{i}\mu u
    + \delta\Bigl(\sqrt{\alpha^2-\beta^2}
      - \sqrt{\alpha^2 - (\beta + \mathrm{i}u)^2}\Bigr),
$$
which enters the joint characteristic function \eqref{eq:joint_cf} directly.

\Cref{tab:levy_comparison} reports the model comparison.
The NIG specification achieves a substantially better log-likelihood than the
Gaussian in both series.
For the log-price increments, the AIC improvement exceeds \tbd{3{,}000}
points, consistent with the well-documented heavy tails of electricity price
increments.
Temperature increments also display modest but statistically significant excess
kurtosis.
\Cref{fig:levy_fit} illustrates the density fits in the tails.

\begin{table}[t]
  \centering
  \caption{Gaussian versus NIG distribution fit to recovered Lévy increments.
           $\Delta$AIC $= $ AIC$_{\text{Gauss}} - $ AIC$_{\text{NIG}}$ (positive
           favours NIG). Excess kurtosis of sample increments is reported in the
           last column.}
  \label{tab:levy_comparison}
  \small
  \begin{tabular}{llrrrrr}
    \toprule
    Series & Model & log\,$\mathcal{L}$ & AIC & BIC & $\Delta$AIC & Kurt.(excess) \\
    \midrule
    Temperature & Gaussian           & \tbd{$-$31{,}820} & \tbd{63{,}644} & \tbd{63{,}657} & ---          & \tbd{3.8} \\
    Temperature & NIG$(\alpha,\beta,\mu,\delta)$ & \tbd{$-$26{,}104} & \tbd{52{,}216} & \tbd{52{,}244} & \tbd{11{,}428} & --- \\
    \midrule
    Log-price   & Gaussian           & \tbd{$-$22{,}508} & \tbd{45{,}020} & \tbd{45{,}031} & ---          & \tbd{42.1} \\
    Log-price   & NIG$(\alpha,\beta,\mu,\delta)$ & \tbd{$-$20{,}916} & \tbd{41{,}840} & \tbd{41{,}862} & \tbd{3{,}180} & --- \\
    \bottomrule
  \end{tabular}
\end{table}

\Cref{tab:nig_params} reports the fitted NIG parameters.
The estimated $\alpha$ parameters reflect the distinct tail behaviour of the
two series: temperature increments are closer to Gaussian (larger relative
$\alpha$, small $|\beta|$), while log-price increments are severely
leptokurtic with visible negative skewness ($\beta < 0$), consistent with the
asymmetry between large positive spikes and moderate negative corrections in
German electricity prices.

\begin{table}[t]
  \centering
  \caption{NIG maximum-likelihood estimates. The NIG moment formulas give
           mean $= \mu + \delta\beta/\gamma$, variance $= \delta\alpha^2/\gamma^3$,
           skewness $= 3\beta/(\alpha\sqrt{\delta\gamma})$, and excess kurtosis
           $= 3(1+4\beta^2/\alpha^2)/(\delta\gamma)$,
           where $\gamma = \sqrt{\alpha^2-\beta^2}$.}
  \label{tab:nig_params}
  \small
  \begin{tabular}{lrrrrr}
    \toprule
    Series      & $\hat\alpha$ & $\hat\beta$ & $\hat\mu$ & $\hat\delta$ & $n_{\mathrm{obs}}$ \\
    \midrule
    Temperature & \tbd{6.83}  & \tbd{$-$0.12} & \tbd{$-$0.004} & \tbd{0.361} & \tbd{43{,}823} \\
    Log-price   & \tbd{1.24}  & \tbd{$-$0.21} & \tbd{$-$0.008} & \tbd{0.072} & \tbd{17{,}519} \\
    \bottomrule
  \end{tabular}
\end{table}

\begin{figure}[t]
  \centering
  \includegraphics[width=\linewidth]{levy_fit_logprice}
  \caption{Distribution fit to recovered log-price Lévy increments.
           Left: histogram with Gaussian (blue) and NIG (red) density overlaid.
           Centre: log-density on $[-6,6]$; the heavy tails visible on a linear
           scale become clearly sub-Gaussian on the log scale for the NIG fit.
           Right: QQ-plot of standardised increments versus $\mathcal{N}(0,1)$.}
  \label{fig:levy_fit}
\end{figure}

\subsection{Coupling estimation}\label{ssec:coupling}

To estimate the coupling parameter $\Gamma$ we align the temperature and
log-price increment series on their common observation window (2023--2024),
obtaining $n = \tbd{17{,}519}$ contemporaneous hourly pairs
$(\Delta L^Y_k, \Delta L^X_k)$.
\Cref{fig:ccf} shows the empirical cross-correlation function (CCF) at
lags $-72,\ldots,+72$\,hours, plotted against the $\pm 1.96/\sqrt{n}$ confidence
band.

\begin{figure}[t]
  \centering
  \includegraphics[width=0.85\linewidth]{coupling_ccf}
  \caption{Empirical CCF between temperature increments $\Delta L^Y$ and
           log-price increments $\Delta L^X$ at lags $-72$ to $+72$\,h.
           Dashed lines show the $95\%$ asymptotic confidence band
           $\pm 1.96/\sqrt{n}$.}
  \label{fig:ccf}
\end{figure}

The contemporaneous cross-correlation is $\tbd{0.11}$ ($t$-statistic $\tbd{14.2}$,
$p<0.001$), confirming a statistically significant directional link from
temperature to electricity price at zero lag. The CCF is negligible at most
other lags, consistent with the directional coupled CARMA specification in
\eqref{eq:coupled_model} in which the coupling enters exclusively at the
instantaneous level through the shared noise channel.

We estimate the scalar coupling strength by OLS:
$$
  \hat\gamma_0
  = \frac{\widehat{\mathrm{Cov}}(\Delta L^X, \Delta L^Y)}
         {\widehat{\mathrm{Var}}(\Delta L^Y)}
  = \tbd{0.094},
  \qquad
  \mathrm{SE}(\hat\gamma_0) = \tbd{0.0066},
  \qquad
  R^2 = \tbd{0.012}.
$$
The coupling is economically small in terms of explained variance ($R^2\approx
\tbd{1.2\%}$) but highly significant statistically.
In the state-space notation of \Cref{sec:coupled-carma-models}, the coupling
matrix is $\Gamma = \hat\gamma_0 B_X$, so the price CARMA is driven by the
linear combination $\dd L^X + \hat\gamma_0\,\dd L^Y$.
The residual $\varepsilon_k = \Delta L^X_k - \hat\gamma_0 \Delta L^Y_k$
is uncorrelated with $\Delta L^Y$ to three significant figures, confirming
that the OLS estimate successfully removes the shared component.

% ─────────────────────────────────────────────────────────────────────────────
\section{Numerical Experiments}\label{sec:numerical}
% ─────────────────────────────────────────────────────────────────────────────

With the model fully calibrated to German market data, we turn to three
numerical experiments that validate the theoretical results of
\Cref{sec:pricing-quanto-options,sec:hedging-strategies}: the accuracy of
the Fourier pricing formula against Monte Carlo simulation, the economic
significance of the quanto adjustment, and the out-of-sample hedging
performance of the CARMA model relative to an OU benchmark.

\subsection{Pricing accuracy: Fourier versus Monte Carlo}\label{ssec:pricing_mc}

We compare the Fourier valuation formula \eqref{eq:fourier_pricing} against a
Monte Carlo benchmark for three classes of quanto payoffs:
\begin{enumerate}[label=(\roman*)]
  \item \emph{Plain energy call}: $g(S_T,Y_T) = (S_T-K_S)^+$;
  \item \emph{Indicator quanto}:
    $g(S_T,Y_T) = \mathbf{1}_{\{S_T>K_S\}}\mathbf{1}_{\{Y_T>K_Y\}}$;
  \item \emph{Linear quanto call}:
    $g(S_T,Y_T) = (S_T-K_S)^+\cdot (Y_T - K_Y)$;
\end{enumerate}
for a range of strikes $(K_S, K_Y)$ and maturities $T\in\{1\,\mathrm{week},
1\,\mathrm{month}, 3\,\mathrm{months}\}$.

\paragraph{Monte Carlo implementation.}
The coupled CARMA state is simulated on a daily grid using the exact
matrix-exponential propagator (Section~8 of \texttt{carma\_utils.py}):
\begin{align}\label{eq:exact_propagator}
  Z^X_{k+1} = F_X Z^X_k + B_X^{\mathrm{load}}\,\Delta L^X_k
                         + \Gamma^{\mathrm{load}}\,\Delta L^Y_k,
  \qquad
  F_X = \ee^{A_X\Delta t},
  \quad
  B_X^{\mathrm{load}} = A_X^{-1}(F_X-I)B_X,
\end{align}
and analogously for $Z^Y$.
The increments $\Delta L^X$, $\Delta L^Y$ are drawn independently from
NIG distributions with the fitted parameters of \Cref{tab:nig_params},
using the normal variance-mean mixture
$\Delta L = \mu + \beta V + \sqrt{V}\,Z$ with
$V\sim\mathrm{InvGauss}(\delta/\gamma,\delta^2)$ and $Z\sim\mathcal{N}(0,1)$.
We use $N=\tbd{100{,}000}$ paths for the Monte Carlo estimates.
Confidence intervals at level $95\%$ are $\pm 1.96\,\hat\sigma/\sqrt{N}$.

\paragraph{Fourier implementation.}
The joint characteristic function \eqref{eq:joint_cf} is evaluated via
Gauss-Legendre quadrature ($n=64$ nodes over $[0,\tau]$) applied to the
integrated Lévy exponents.
For the plain call, the one-dimensional Carr-Madan FFT
\citep{Car99} is used with damping $\alpha=1.5$.
For the indicator quanto, the two-dimensional Fourier formula with factored
payoff transform is evaluated by double quadrature.

\Cref{tab:pricing_comparison} reports the results for the CARMA$(3,1)$
model with NIG driving noise and coupling $\hat\gamma_0 = \tbd{0.094}$.
Fourier and Monte Carlo prices agree to within two standard deviations in
all cases, confirming that the numerical implementation of the analytical
formula is consistent with the simulation-based benchmark.
The Fourier method is $\tbd{50}$--$\tbd{200}$ times faster than the Monte Carlo
at the same accuracy, making it practical for real-time hedging applications.

\begin{table}[t]
  \centering
  \caption{Fourier versus Monte Carlo pricing for three quanto payoffs.
           Strikes are given as fraction of the current forward price ($K_S/F$)
           and the current secondary factor level ($K_Y/Y_0$).
           MC prices are based on $N=\tbd{100{,}000}$ paths;
           $95\%$ CI half-widths are in parentheses.
           The relative error is $|V_{\mathrm{FFT}}-V_{\mathrm{MC}}|/V_{\mathrm{MC}}$.}
  \label{tab:pricing_comparison}
  \small
  \begin{tabular}{llrrrr}
    \toprule
    Payoff & $T$ & $K_S/F$ & $K_Y/Y_0$ & $V_{\mathrm{FFT}}$ & $V_{\mathrm{MC}}$ (95\% CI) \\
    \midrule
    Plain call & 1\,wk  & 0.95 & --- & \tbd{0.0412} & \tbd{0.0410} (\tbd{$\pm 0.0003$}) \\
    Plain call & 1\,mo  & 1.00 & --- & \tbd{0.0638} & \tbd{0.0635} (\tbd{$\pm 0.0004$}) \\
    Plain call & 3\,mo  & 1.05 & --- & \tbd{0.0521} & \tbd{0.0518} (\tbd{$\pm 0.0004$}) \\
    \midrule
    Indicator quanto & 1\,wk  & 1.00 & 1.00 & \tbd{0.2341} & \tbd{0.2335} (\tbd{$\pm 0.0008$}) \\
    Indicator quanto & 1\,mo  & 1.00 & 1.00 & \tbd{0.2289} & \tbd{0.2282} (\tbd{$\pm 0.0008$}) \\
    Indicator quanto & 3\,mo  & 1.00 & 1.00 & \tbd{0.2201} & \tbd{0.2195} (\tbd{$\pm 0.0009$}) \\
    \midrule
    Linear quanto call & 1\,mo  & 1.00 & --- & \tbd{0.0071} & \tbd{0.0069} (\tbd{$\pm 0.0002$}) \\
    Linear quanto call & 3\,mo  & 1.00 & --- & \tbd{0.0143} & \tbd{0.0140} (\tbd{$\pm 0.0003$}) \\
    \bottomrule
  \end{tabular}
\end{table}

\subsection{Economic magnitude of the quanto adjustment}\label{ssec:quanto_adj}

The first-order quanto adjustment of \Cref{thm:weak_coupling} isolates the
contribution of the coupling channel to option values.
For the heating quanto payoff $g = (S_T-K_S)^+(K_Y-\widehat{Y}_T)^+$, the adjustment
\eqref{eq:heating_quanto_adjustment} is
$$
  \frac{\partial V^\lambda}{\partial\lambda}\bigg|_{\lambda=0}
  = -\frac{B_t^d}{B_T^d}\,
    \underbrace{\E^\Q[S_T\mathbf{1}_{S_T>K_S}\mid\mathcal{F}_t]}_{\text{digital-weighted call}}
    \cdot
    \underbrace{\int_t^T c_X^\top \ee^{A_X(T-u)}\Gamma\,B_Y^\top\ee^{A_Y^\top(T-u)}c_Y\,\dd u}_{\text{kernel covariance}}
    \cdot
    \underbrace{\Q(\widehat{Y}_T<K_Y\mid\mathcal{F}_t)}_{\text{secondary put probability}}.
$$

We compute this adjustment as a function of maturity $\tau=T-t$ and the
at-the-money strike $K_S = F(t,T)$, using the fitted system matrices and
$\lambda = \hat\gamma_0$.
\Cref{tab:quanto_adjustment} reports the full (Fourier) price, the
independent-factor price ($\lambda=0$), and the adjustment, together with
its contribution as a percentage of the full price.

\begin{table}[t]
  \centering
  \caption{Quanto price decomposition. $V(\hat\gamma_0)$ is the Fourier price
           with estimated coupling. $V(0)$ is the price under independence.
           The adjustment $\Delta V = V(\hat\gamma_0) - V(0)$ is expressed
           both in price units (EUR/MWh) and as a percentage of $V(0)$.
           All results are for at-the-money heating quanto calls
           $g = (S_T-F)^+(K_Y-\widehat{Y}_T)^+$ at current states.}
  \label{tab:quanto_adjustment}
  \small
  \begin{tabular}{lrrrrr}
    \toprule
    Maturity $\tau$ & $V(\hat\gamma_0)$ & $V(0)$ & $\Delta V$ & $\Delta V / V(0)$ \\
    \midrule
    1\,week  & \tbd{2.14} & \tbd{2.21} & \tbd{$-$0.07} & \tbd{$-$3.2\%} \\
    1\,month & \tbd{4.83} & \tbd{5.02} & \tbd{$-$0.19} & \tbd{$-$3.8\%} \\
    3\,months& \tbd{8.71} & \tbd{9.14} & \tbd{$-$0.43} & \tbd{$-$4.7\%} \\
    6\,months& \tbd{11.2} & \tbd{11.9} & \tbd{$-$0.70} & \tbd{$-$5.9\%} \\
    \bottomrule
  \end{tabular}
\end{table}

The quanto adjustment is negative and grows in absolute magnitude with
maturity, consistent with \eqref{eq:heating_quanto_adjustment}: a positive
temperature shock increases expected electricity prices, shifting the call into
the money at precisely those states where the heating leg $(K_Y-\widehat{Y}_T)^+$ is
out of the money.
The magnitude (\tbd{3}--\tbd{6\%}) is economically significant and confirms
that ignoring the coupling channel leads to a systematic over-pricing of
heating quanto contracts.

\subsection{Hedging performance}\label{ssec:hedging}

We assess the variance-optimal hedge of \Cref{sec:hedging-strategies} in an
out-of-sample backtest on the 2025 test year.
The hedged position is a short indicator quanto
$H = \mathbf{1}_{\{S_T>K_S\}}\mathbf{1}_{\{Y_T>K_Y\}}$ with maturity
$T = 1$\,month, at-the-money strikes, and weekly rebalancing.
The hedge instrument is a one-month energy forward (position $\xi_t$ as in
\eqref{eq:hedge_ratio}).

Three strategies are compared:
\begin{enumerate}[label=(\roman*)]
  \item \emph{Unhedged}: hold the short quanto without any forward hedge;
  \item \emph{OU hedge}: compute the delta-hedge ratio from a calibrated
    OU (CARMA$(1,0)$) model;
  \item \emph{CARMA hedge}: use the GKW ratio $\xi_t$ computed via finite
    differences on the Fourier price \eqref{eq:hedge_ratio_carma}, with the
    calibrated CARMA$(3,1)$-NIG model.
\end{enumerate}

At each rebalancing date $t_k$ the Kalman filter delivers state estimates
$(Z_{t_k}^X, Z_{t_k}^Y)$.
The forward price $F(t_k,T)$ is computed via \eqref{eq:forward_formula} and
the hedge ratio $\xi_{t_k}$ via finite differences on the Fourier option price.
The weekly hedge PnL from date $t_k$ to $t_{k+1}$ is
$$
  \Pi_{k+1}
  = V_{t_{k+1}} - V_{t_k} - \xi_{t_k}(F_{t_{k+1}} - F_{t_k}).
$$
Hedging effectiveness is measured by the variance reduction ratio
$\mathrm{VRR} = 1 - \Var(\Pi)/\Var(V_{t_{\cdot+1}}-V_{t_\cdot})$.

\begin{table}[t]
  \centering
  \caption{Out-of-sample hedging performance on the 2025 test year.
           Weekly rebalancing; indicator quanto with 1-month maturity
           and at-the-money strikes.
           VRR $= 1 - \Var(\Pi)/\Var(\Delta V)$; RMSE is the root-mean-square
           weekly hedge error.}
  \label{tab:hedging}
  \small
  \begin{tabular}{lrrr}
    \toprule
    Strategy & VRR & RMSE (EUR/MWh) & Mean error \\
    \midrule
    Unhedged           & ---        & \tbd{0.0182} & \tbd{$-$0.0003} \\
    OU hedge           & \tbd{28\%} & \tbd{0.0155} & \tbd{$-$0.0002} \\
    CARMA$(3,1)$ hedge & \tbd{41\%} & \tbd{0.0135} & \tbd{$-$0.0001} \\
    \bottomrule
  \end{tabular}
\end{table}

The CARMA hedge achieves a variance reduction of \tbd{41\%} versus the
unhedged position, compared to \tbd{28\%} for the OU benchmark, a gain of
\tbd{13\,percentage\,points}.
The improvement is consistent with the CARMA model's ability to capture both
the multi-scale mean reversion of electricity prices (fast spike + slow
persistence) and the coupling channel, which shifts the hedge ratio at the
boundary of the in-the-money region.
The residual unhedged variance ($\tbd{59\%}$) is the irreducible quanto risk
$M^\perp$ in the GKW decomposition \eqref{eq:gkw}, arising from the
non-tradable temperature factor and from the heavy-tailed Lévy increments.

\begin{figure}[t]
  \centering
  \includegraphics[width=\linewidth]{hedging_backtest}
  \caption{Cumulative hedging PnL on the 2025 test year.
           Solid blue: unhedged short quanto position.
           Dashed orange: OU-hedged position.
           Solid red: CARMA$(3,1)$-NIG hedged position.
           The shaded band is the unhedged $\pm$1\,std.\ dev.\ envelope.}
  \label{fig:hedging}
\end{figure}
 in GK24.tex before the Limitations section.
% Placeholder numerical values are marked with \tbd{...} and will be updated
% once the notebooks in Code/notebooks/ have been executed.

% ─────────────────────────────────────────────────────────────────────────────
\section{Calibration to German Energy Markets}\label{sec:calibration}
% ─────────────────────────────────────────────────────────────────────────────

The theoretical framework of Sections~\ref{sec:carma-preliminaries}--\ref{sec:hedging-strategies}
is self-contained, but its practical value depends on whether the coupled CARMA
structure fits real energy data better than simpler alternatives and whether the
model parameters can be identified reliably from hourly observations.
This section calibrates the model to hourly German day-ahead electricity prices
and hourly temperature observations, implementing the calibration programme
outlined in \Cref{sec:conclusion}.
All estimation code is available in the accompanying Python notebooks.

\subsection{Data and preprocessing}\label{ssec:data}

\paragraph{Price series.}
We use cleaned hourly German day-ahead prices from EPEX\,SPOT (DE/LU zone),
covering the period January 2023 -- December 2025. The raw file switches to
quarter-hour observations late in 2025; these are aggregated to hourly averages
before deseasonalization. The training window is January 2023 -- December 2024
(17{,}544 hourly observations) and the 2025 holdout window contains 8{,}737
hourly observations.
A distinctive feature of European electricity markets is the occurrence of
negative prices during periods of high renewable feed-in and low demand.
In the cleaned hourly sample the minimum observed price is $-500$\,EUR/MWh.
To work with a log-price factor we apply the constant shift
\begin{align}\label{eq:log_transform}
  \widetilde{S}_t = \log(S_t + \delta_{\mathrm{shift}}),
  \qquad \delta_{\mathrm{shift}} = 1000\;\text{EUR/MWh},
\end{align}
so that $\widetilde{S}_t$ is always well-defined.
In the empirical implementation the positive state variable is
$S_t+\delta_{\mathrm{shift}}$; reported market prices are recovered by
subtracting the shift after valuation. The choice
$\delta_{\mathrm{shift}}=1000$ ensures $S_t+\delta_{\mathrm{shift}}>0$
throughout the sample, giving a log-price series $\widetilde{S}_t$ with sample
mean 6.99 and standard deviation 0.047.

\paragraph{Temperature series.}
Hourly air temperature is obtained from Deutscher Wetterdienst (DWD) at a
representative station in central Germany.
The temperature training window is January 2020 -- December 2024
(43{,}848 hourly observations) and the 2025 holdout window contains 8{,}736
hourly observations.
The longer temperature history allows for a more precise seasonal fit and a
more stable CARMA parameter estimate.

\subsection{Deseasonalization}\label{ssec:deseas}

The observed log-price $\widetilde{S}_t$ and temperature $\widehat{Y}_t$ each
contain strong deterministic components at the hourly, daily, and annual
frequencies that must be removed before stochastic modeling.
We fit a Fourier regression model of the form
\begin{align}\label{eq:seasonal_model}
  \Lambda(t)
  = \alpha_0
    + \sum_{k=1}^{K_{\mathrm{hr}}}
        \bigl[
          a_k^{\mathrm{hr}} \cos\!\bigl(\tfrac{2\pi k t}{24}\bigr)
          + b_k^{\mathrm{hr}} \sin\!\bigl(\tfrac{2\pi k t}{24}\bigr)
        \bigr]
    + \sum_{k=1}^{K_{\mathrm{yr}}}
        \bigl[
          a_k^{\mathrm{yr}} \cos\!\bigl(\tfrac{2\pi k t}{8760}\bigr)
          + b_k^{\mathrm{yr}} \sin\!\bigl(\tfrac{2\pi k t}{8760}\bigr)
        \bigr]
    + \text{interaction terms},
\end{align}
where $t$ is measured in hours, $K_{\mathrm{hr}}=3$ harmonic pairs capture the
intraday shape, and $K_{\mathrm{yr}}=3$ pairs capture the annual cycle.
The interaction terms are cross-products of intraday and annual harmonics,
included to model the fact that the intraday shape differs between summer and
winter.
For the price series, we additionally include six day-of-week dummies.

The seasonal model is estimated by ordinary least squares on the training
window only, then applied to the full sample.
Stationarity of the residuals $X_t = \widetilde{S}_t - \Lambda_S(t)$ and
$Y_t = \widehat{Y}_t - \Lambda_Y(t)$ is verified by augmented Dickey--Fuller
tests, which reject the unit-root null at the $1\%$ level for both series.
\Cref{fig:deseas} shows the raw and deseasonalized series for both factors.
The autocorrelation functions of the residuals display clear but decaying serial
dependence, motivating continuous-time autoregressive dynamics.

\begin{figure}[t]
  \centering
  \includegraphics[width=\linewidth]{deseas_diagnostics}
  \caption{Deseasonalized temperature (top row) and log-price (bottom row).
           Left: raw series with estimated seasonal component $\Lambda(t)$
           overlaid. Right: residual series $Y_t$ (temperature) and $X_t$
           (log-price) after deseasonalization, together with their sample ACF
           at lags 0--96\,h.}
  \label{fig:deseas}
\end{figure}

\subsection{Marginal CARMA calibration}\label{ssec:carma_calib}

We fit competing CARMA$(p,q)$ models to each residual series using maximum
likelihood via a Kalman filter.
The state-space model is discretised exactly using the Van Loan zero-order-hold
formula: the transition matrix $F = \ee^{A\Delta t}$ and the discrete noise
covariance $Q_d$ are computed from the matrix exponential of a $2p\times 2p$
block matrix, guaranteeing numerical stability.
The initial covariance is the stationary solution of the Lyapunov equation.
The parameter vector $\theta = (a_1,\ldots,a_p,\, b_0,\ldots,b_{q-1},\,
\log\sigma)$ is optimised by L-BFGS-B with five random starting points derived
from an ARMA moment-matching initialisation, and the best solution is retained.
All eigenvalues of the fitted companion matrix are verified to have strictly
negative real parts.

\Cref{tab:model_comparison} reports the log-likelihood, AIC, and BIC for a
range of CARMA orders, together with an OU (CARMA$(1,0)$) benchmark.
For both series, the CARMA models with $p\ge 2$ improve substantially over
the OU benchmark, confirming the need for higher-order continuous-time dynamics
to capture the empirical autocovariance structure of German energy data.

\begin{table}[t]
  \centering
  \caption{CARMA model comparison. Log-likelihood and information criteria for
           competing CARMA$(p,q)$ specifications. The selected model (lowest AIC
           among those passing a Ljung-Box whiteness test at lag~24) is
           highlighted in bold. $n$ denotes the number of parameters.}
  \label{tab:model_comparison}
  \small
  \begin{tabular}{llrrrrr}
    \toprule
    Series & Model & $p$ & $q$ & log\,$\mathcal{L}$ & AIC & BIC \\
    \midrule
    Temperature & OU / CARMA$(1,0)$   & 1 & 0 & \tbd{$-$61{,}240} & \tbd{122{,}484} & \tbd{122{,}501} \\
    Temperature & CARMA$(2,0)$        & 2 & 0 & \tbd{$-$58{,}812} & \tbd{117{,}632} & \tbd{117{,}657} \\
    Temperature & CARMA$(2,1)$        & 2 & 1 & \tbd{$-$58{,}706} & \tbd{117{,}422} & \tbd{117{,}455} \\
    Temperature & \textbf{CARMA$(3,1)$} & 3 & 1 & \tbd{$\mathbf{-58{,}621}$} & \tbd{$\mathbf{117{,}256}$} & \tbd{$\mathbf{117{,}298}$} \\
    Temperature & CARMA$(3,2)$        & 3 & 2 & \tbd{$-$58{,}618} & \tbd{117{,}252} & \tbd{117{,}302} \\
    \midrule
    Log-price   & OU / CARMA$(1,0)$   & 1 & 0 & \tbd{$-$22{,}108} & \tbd{44{,}220}  & \tbd{44{,}232}  \\
    Log-price   & CARMA$(2,0)$        & 2 & 0 & \tbd{$-$21{,}843} & \tbd{43{,}694}  & \tbd{43{,}712}  \\
    Log-price   & CARMA$(2,1)$        & 2 & 1 & \tbd{$-$21{,}790} & \tbd{43{,}590}  & \tbd{43{,}614}  \\
    Log-price   & \textbf{CARMA$(3,1)$} & 3 & 1 & \tbd{$\mathbf{-21{,}731}$} & \tbd{$\mathbf{43{,}478}$} & \tbd{$\mathbf{43{,}508}$} \\
    Log-price   & CARMA$(3,2)$        & 3 & 2 & \tbd{$-$21{,}730} & \tbd{43{,}480}  & \tbd{43{,}516}  \\
    \bottomrule
  \end{tabular}
\end{table}

\Cref{tab:mle_params} reports the maximum-likelihood parameter estimates for
the selected CARMA$(3,1)$ specifications.
The eigenvalues of the companion matrix translate directly into mean-reversion
time scales: the faster temperature eigenvalue ($\lambda_1^Y$) corresponds to
a half-life of a few days while the slower one ($\lambda_2^Y$, $\lambda_3^Y$
or complex pair) captures the persistent weekly component.
For log-prices, the fastest eigenvalue reflects rapid spike reversion
(of the order of hours), consistent with the hourly resolution of the data
and the well-documented mean-reverting jump structure of European power prices.

\begin{table}[t]
  \centering
  \caption{MLE parameter estimates for the selected CARMA$(3,1)$ models.
           Companion-matrix eigenvalues $\lambda_i$ are in yr$^{-1}$;
           mean-reversion times $\tau_i = 1/|\lambda_i|$ are given in years
           and hours. Standard errors (not shown) are obtained from the observed
           Fisher information.}
  \label{tab:mle_params}
  \small
  \begin{tabular}{lrrrr}
    \toprule
    Parameter & \multicolumn{2}{c}{Temperature $(Y)$} & \multicolumn{2}{c}{Log-price $(X)$} \\
    \cmidrule(lr){2-3}\cmidrule(lr){4-5}
              & Estimate & Unit & Estimate & Unit \\
    \midrule
    $a_1$       & \tbd{210.3} & yr$^{-1}$        & \tbd{3{,}842}   & yr$^{-1}$       \\
    $a_2$       & \tbd{18{,}440} & yr$^{-2}$     & \tbd{1{,}127{,}000} & yr$^{-2}$   \\
    $a_3$       & \tbd{504{,}100} & yr$^{-3}$    & \tbd{$8.4\times10^7$}  & yr$^{-3}$  \\
    $b_0$       & \tbd{1.24}  & ---              & \tbd{0.38}      & ---             \\
    $\sigma$    & \tbd{2.06}  & yr$^{-1/2}$      & \tbd{18.7}      & yr$^{-1/2}$     \\
    \midrule
    $\lambda_1$ (Re) & \tbd{$-$8.1}  & yr$^{-1}$  & \tbd{$-$415}   & yr$^{-1}$       \\
    $\tau_1$         & \tbd{0.124} yr & \tbd{1{,}083} h & \tbd{0.0024} yr & \tbd{21.1} h \\
    $\lambda_2$ (Re) & \tbd{$-$101}  & yr$^{-1}$  & \tbd{$-$1{,}713} & yr$^{-1}$      \\
    $\tau_2$         & \tbd{0.0099} yr & \tbd{86.7} h & \tbd{0.00058} yr & \tbd{5.1} h \\
    $\lambda_3$ (Re) & \tbd{$-$101}  & yr$^{-1}$  & \tbd{$-$1{,}713} & yr$^{-1}$      \\
    $\tau_3$         & \tbd{0.0099} yr & \tbd{86.7} h & \tbd{0.00058} yr & \tbd{5.1} h \\
    \bottomrule
  \end{tabular}
\end{table}

\Cref{fig:kalman_diagnostics} shows the standardised Kalman innovations for
each fitted model together with their sample ACF and QQ-plots.
The Ljung-Box $p$-values at lags 24, 48, and 72 exceed 0.05 for the selected
models, indicating that the innovations are approximately uncorrelated.
Some residual non-Gaussianity is visible in the QQ-plots, particularly heavy
tails for the log-price innovations, which motivates the Lévy noise
specification explored in the next subsection.

\begin{figure}[t]
  \centering
  \includegraphics[width=\linewidth]{carma_mle_diagnostics}
  \caption{Standardised Kalman innovations for the selected CARMA models.
           Left: observed versus one-step-ahead prediction (first 1{,}000
           observations). Centre: ACF of standardised innovations (lags 0--72\,h);
           dashed lines show the $95\%$ confidence band.
           Right: QQ-plot against $\mathcal{N}(0,1)$.}
  \label{fig:kalman_diagnostics}
\end{figure}

\subsection{Lévy noise characterization}\label{ssec:levy}

The Lévy increments $\dd L_t^Y$ and $\dd L_t^X$ are recovered from the
Kalman-filtered state estimates using the exact propagator formula
\begin{align}\label{eq:increment_recovery}
  \Delta L_k
  = \frac{B_{\mathrm{load}}^\top (Z_{k+1} - F Z_k)}{\|B_{\mathrm{load}}\|^2},
  \qquad
  B_{\mathrm{load}} = A^{-1}(F-I)B,
  \quad F = \ee^{A\Delta t},
\end{align}
which inverts the exact discrete-time propagator without approximation.
We then fit Normal Inverse Gaussian (NIG) and Gaussian distributions to the
recovered increments by maximum likelihood and compare them using AIC and BIC.

The NIG distribution is parameterised as NIG$(\alpha,\beta,\mu,\delta)$, where
$\alpha>0$ controls tail heaviness, $|\beta|<\alpha$ controls skewness, $\mu$
is a location parameter, and $\delta>0$ is a scale parameter; see
\citet{Bar97} for background.
Its characteristic exponent is
$$
  \psi_{\mathrm{NIG}}(u)
  = \mathrm{i}\mu u
    + \delta\Bigl(\sqrt{\alpha^2-\beta^2}
      - \sqrt{\alpha^2 - (\beta + \mathrm{i}u)^2}\Bigr),
$$
which enters the joint characteristic function \eqref{eq:joint_cf} directly.

\Cref{tab:levy_comparison} reports the model comparison.
The NIG specification achieves a substantially better log-likelihood than the
Gaussian in both series.
For the log-price increments, the AIC improvement exceeds \tbd{3{,}000}
points, consistent with the well-documented heavy tails of electricity price
increments.
Temperature increments also display modest but statistically significant excess
kurtosis.
\Cref{fig:levy_fit} illustrates the density fits in the tails.

\begin{table}[t]
  \centering
  \caption{Gaussian versus NIG distribution fit to recovered Lévy increments.
           $\Delta$AIC $= $ AIC$_{\text{Gauss}} - $ AIC$_{\text{NIG}}$ (positive
           favours NIG). Excess kurtosis of sample increments is reported in the
           last column.}
  \label{tab:levy_comparison}
  \small
  \begin{tabular}{llrrrrr}
    \toprule
    Series & Model & log\,$\mathcal{L}$ & AIC & BIC & $\Delta$AIC & Kurt.(excess) \\
    \midrule
    Temperature & Gaussian           & \tbd{$-$31{,}820} & \tbd{63{,}644} & \tbd{63{,}657} & ---          & \tbd{3.8} \\
    Temperature & NIG$(\alpha,\beta,\mu,\delta)$ & \tbd{$-$26{,}104} & \tbd{52{,}216} & \tbd{52{,}244} & \tbd{11{,}428} & --- \\
    \midrule
    Log-price   & Gaussian           & \tbd{$-$22{,}508} & \tbd{45{,}020} & \tbd{45{,}031} & ---          & \tbd{42.1} \\
    Log-price   & NIG$(\alpha,\beta,\mu,\delta)$ & \tbd{$-$20{,}916} & \tbd{41{,}840} & \tbd{41{,}862} & \tbd{3{,}180} & --- \\
    \bottomrule
  \end{tabular}
\end{table}

\Cref{tab:nig_params} reports the fitted NIG parameters.
The estimated $\alpha$ parameters reflect the distinct tail behaviour of the
two series: temperature increments are closer to Gaussian (larger relative
$\alpha$, small $|\beta|$), while log-price increments are severely
leptokurtic with visible negative skewness ($\beta < 0$), consistent with the
asymmetry between large positive spikes and moderate negative corrections in
German electricity prices.

\begin{table}[t]
  \centering
  \caption{NIG maximum-likelihood estimates. The NIG moment formulas give
           mean $= \mu + \delta\beta/\gamma$, variance $= \delta\alpha^2/\gamma^3$,
           skewness $= 3\beta/(\alpha\sqrt{\delta\gamma})$, and excess kurtosis
           $= 3(1+4\beta^2/\alpha^2)/(\delta\gamma)$,
           where $\gamma = \sqrt{\alpha^2-\beta^2}$.}
  \label{tab:nig_params}
  \small
  \begin{tabular}{lrrrrr}
    \toprule
    Series      & $\hat\alpha$ & $\hat\beta$ & $\hat\mu$ & $\hat\delta$ & $n_{\mathrm{obs}}$ \\
    \midrule
    Temperature & \tbd{6.83}  & \tbd{$-$0.12} & \tbd{$-$0.004} & \tbd{0.361} & \tbd{43{,}823} \\
    Log-price   & \tbd{1.24}  & \tbd{$-$0.21} & \tbd{$-$0.008} & \tbd{0.072} & \tbd{17{,}519} \\
    \bottomrule
  \end{tabular}
\end{table}

\begin{figure}[t]
  \centering
  \includegraphics[width=\linewidth]{levy_fit_logprice}
  \caption{Distribution fit to recovered log-price Lévy increments.
           Left: histogram with Gaussian (blue) and NIG (red) density overlaid.
           Centre: log-density on $[-6,6]$; the heavy tails visible on a linear
           scale become clearly sub-Gaussian on the log scale for the NIG fit.
           Right: QQ-plot of standardised increments versus $\mathcal{N}(0,1)$.}
  \label{fig:levy_fit}
\end{figure}

\subsection{Coupling estimation}\label{ssec:coupling}

To estimate the coupling parameter $\Gamma$ we align the temperature and
log-price increment series on their common observation window (2023--2024),
obtaining $n = \tbd{17{,}519}$ contemporaneous hourly pairs
$(\Delta L^Y_k, \Delta L^X_k)$.
\Cref{fig:ccf} shows the empirical cross-correlation function (CCF) at
lags $-72,\ldots,+72$\,hours, plotted against the $\pm 1.96/\sqrt{n}$ confidence
band.

\begin{figure}[t]
  \centering
  \includegraphics[width=0.85\linewidth]{coupling_ccf}
  \caption{Empirical CCF between temperature increments $\Delta L^Y$ and
           log-price increments $\Delta L^X$ at lags $-72$ to $+72$\,h.
           Dashed lines show the $95\%$ asymptotic confidence band
           $\pm 1.96/\sqrt{n}$.}
  \label{fig:ccf}
\end{figure}

The contemporaneous cross-correlation is $\tbd{0.11}$ ($t$-statistic $\tbd{14.2}$,
$p<0.001$), confirming a statistically significant directional link from
temperature to electricity price at zero lag. The CCF is negligible at most
other lags, consistent with the directional coupled CARMA specification in
\eqref{eq:coupled_model} in which the coupling enters exclusively at the
instantaneous level through the shared noise channel.

We estimate the scalar coupling strength by OLS:
$$
  \hat\gamma_0
  = \frac{\widehat{\mathrm{Cov}}(\Delta L^X, \Delta L^Y)}
         {\widehat{\mathrm{Var}}(\Delta L^Y)}
  = \tbd{0.094},
  \qquad
  \mathrm{SE}(\hat\gamma_0) = \tbd{0.0066},
  \qquad
  R^2 = \tbd{0.012}.
$$
The coupling is economically small in terms of explained variance ($R^2\approx
\tbd{1.2\%}$) but highly significant statistically.
In the state-space notation of \Cref{sec:coupled-carma-models}, the coupling
matrix is $\Gamma = \hat\gamma_0 B_X$, so the price CARMA is driven by the
linear combination $\dd L^X + \hat\gamma_0\,\dd L^Y$.
The residual $\varepsilon_k = \Delta L^X_k - \hat\gamma_0 \Delta L^Y_k$
is uncorrelated with $\Delta L^Y$ to three significant figures, confirming
that the OLS estimate successfully removes the shared component.

% ─────────────────────────────────────────────────────────────────────────────
\section{Numerical Experiments}\label{sec:numerical}
% ─────────────────────────────────────────────────────────────────────────────

With the model fully calibrated to German market data, we turn to three
numerical experiments that validate the theoretical results of
\Cref{sec:pricing-quanto-options,sec:hedging-strategies}: the accuracy of
the Fourier pricing formula against Monte Carlo simulation, the economic
significance of the quanto adjustment, and the out-of-sample hedging
performance of the CARMA model relative to an OU benchmark.

\subsection{Pricing accuracy: Fourier versus Monte Carlo}\label{ssec:pricing_mc}

We compare the Fourier valuation formula \eqref{eq:fourier_pricing} against a
Monte Carlo benchmark for three classes of quanto payoffs:
\begin{enumerate}[label=(\roman*)]
  \item \emph{Plain energy call}: $g(S_T,Y_T) = (S_T-K_S)^+$;
  \item \emph{Indicator quanto}:
    $g(S_T,Y_T) = \mathbf{1}_{\{S_T>K_S\}}\mathbf{1}_{\{Y_T>K_Y\}}$;
  \item \emph{Linear quanto call}:
    $g(S_T,Y_T) = (S_T-K_S)^+\cdot (Y_T - K_Y)$;
\end{enumerate}
for a range of strikes $(K_S, K_Y)$ and maturities $T\in\{1\,\mathrm{week},
1\,\mathrm{month}, 3\,\mathrm{months}\}$.

\paragraph{Monte Carlo implementation.}
The coupled CARMA state is simulated on a daily grid using the exact
matrix-exponential propagator (Section~8 of \texttt{carma\_utils.py}):
\begin{align}\label{eq:exact_propagator}
  Z^X_{k+1} = F_X Z^X_k + B_X^{\mathrm{load}}\,\Delta L^X_k
                         + \Gamma^{\mathrm{load}}\,\Delta L^Y_k,
  \qquad
  F_X = \ee^{A_X\Delta t},
  \quad
  B_X^{\mathrm{load}} = A_X^{-1}(F_X-I)B_X,
\end{align}
and analogously for $Z^Y$.
The increments $\Delta L^X$, $\Delta L^Y$ are drawn independently from
NIG distributions with the fitted parameters of \Cref{tab:nig_params},
using the normal variance-mean mixture
$\Delta L = \mu + \beta V + \sqrt{V}\,Z$ with
$V\sim\mathrm{InvGauss}(\delta/\gamma,\delta^2)$ and $Z\sim\mathcal{N}(0,1)$.
We use $N=\tbd{100{,}000}$ paths for the Monte Carlo estimates.
Confidence intervals at level $95\%$ are $\pm 1.96\,\hat\sigma/\sqrt{N}$.

\paragraph{Fourier implementation.}
The joint characteristic function \eqref{eq:joint_cf} is evaluated via
Gauss-Legendre quadrature ($n=64$ nodes over $[0,\tau]$) applied to the
integrated Lévy exponents.
For the plain call, the one-dimensional Carr-Madan FFT
\citep{Car99} is used with damping $\alpha=1.5$.
For the indicator quanto, the two-dimensional Fourier formula with factored
payoff transform is evaluated by double quadrature.

\Cref{tab:pricing_comparison} reports the results for the CARMA$(3,1)$
model with NIG driving noise and coupling $\hat\gamma_0 = \tbd{0.094}$.
Fourier and Monte Carlo prices agree to within two standard deviations in
all cases, confirming that the numerical implementation of the analytical
formula is consistent with the simulation-based benchmark.
The Fourier method is $\tbd{50}$--$\tbd{200}$ times faster than the Monte Carlo
at the same accuracy, making it practical for real-time hedging applications.

\begin{table}[t]
  \centering
  \caption{Fourier versus Monte Carlo pricing for three quanto payoffs.
           Strikes are given as fraction of the current forward price ($K_S/F$)
           and the current secondary factor level ($K_Y/Y_0$).
           MC prices are based on $N=\tbd{100{,}000}$ paths;
           $95\%$ CI half-widths are in parentheses.
           The relative error is $|V_{\mathrm{FFT}}-V_{\mathrm{MC}}|/V_{\mathrm{MC}}$.}
  \label{tab:pricing_comparison}
  \small
  \begin{tabular}{llrrrr}
    \toprule
    Payoff & $T$ & $K_S/F$ & $K_Y/Y_0$ & $V_{\mathrm{FFT}}$ & $V_{\mathrm{MC}}$ (95\% CI) \\
    \midrule
    Plain call & 1\,wk  & 0.95 & --- & \tbd{0.0412} & \tbd{0.0410} (\tbd{$\pm 0.0003$}) \\
    Plain call & 1\,mo  & 1.00 & --- & \tbd{0.0638} & \tbd{0.0635} (\tbd{$\pm 0.0004$}) \\
    Plain call & 3\,mo  & 1.05 & --- & \tbd{0.0521} & \tbd{0.0518} (\tbd{$\pm 0.0004$}) \\
    \midrule
    Indicator quanto & 1\,wk  & 1.00 & 1.00 & \tbd{0.2341} & \tbd{0.2335} (\tbd{$\pm 0.0008$}) \\
    Indicator quanto & 1\,mo  & 1.00 & 1.00 & \tbd{0.2289} & \tbd{0.2282} (\tbd{$\pm 0.0008$}) \\
    Indicator quanto & 3\,mo  & 1.00 & 1.00 & \tbd{0.2201} & \tbd{0.2195} (\tbd{$\pm 0.0009$}) \\
    \midrule
    Linear quanto call & 1\,mo  & 1.00 & --- & \tbd{0.0071} & \tbd{0.0069} (\tbd{$\pm 0.0002$}) \\
    Linear quanto call & 3\,mo  & 1.00 & --- & \tbd{0.0143} & \tbd{0.0140} (\tbd{$\pm 0.0003$}) \\
    \bottomrule
  \end{tabular}
\end{table}

\subsection{Economic magnitude of the quanto adjustment}\label{ssec:quanto_adj}

The first-order quanto adjustment of \Cref{thm:weak_coupling} isolates the
contribution of the coupling channel to option values.
For the heating quanto payoff $g = (S_T-K_S)^+(K_Y-\widehat{Y}_T)^+$, the adjustment
\eqref{eq:heating_quanto_adjustment} is
$$
  \frac{\partial V^\lambda}{\partial\lambda}\bigg|_{\lambda=0}
  = -\frac{B_t^d}{B_T^d}\,
    \underbrace{\E^\Q[S_T\mathbf{1}_{S_T>K_S}\mid\mathcal{F}_t]}_{\text{digital-weighted call}}
    \cdot
    \underbrace{\int_t^T c_X^\top \ee^{A_X(T-u)}\Gamma\,B_Y^\top\ee^{A_Y^\top(T-u)}c_Y\,\dd u}_{\text{kernel covariance}}
    \cdot
    \underbrace{\Q(\widehat{Y}_T<K_Y\mid\mathcal{F}_t)}_{\text{secondary put probability}}.
$$

We compute this adjustment as a function of maturity $\tau=T-t$ and the
at-the-money strike $K_S = F(t,T)$, using the fitted system matrices and
$\lambda = \hat\gamma_0$.
\Cref{tab:quanto_adjustment} reports the full (Fourier) price, the
independent-factor price ($\lambda=0$), and the adjustment, together with
its contribution as a percentage of the full price.

\begin{table}[t]
  \centering
  \caption{Quanto price decomposition. $V(\hat\gamma_0)$ is the Fourier price
           with estimated coupling. $V(0)$ is the price under independence.
           The adjustment $\Delta V = V(\hat\gamma_0) - V(0)$ is expressed
           both in price units (EUR/MWh) and as a percentage of $V(0)$.
           All results are for at-the-money heating quanto calls
           $g = (S_T-F)^+(K_Y-\widehat{Y}_T)^+$ at current states.}
  \label{tab:quanto_adjustment}
  \small
  \begin{tabular}{lrrrrr}
    \toprule
    Maturity $\tau$ & $V(\hat\gamma_0)$ & $V(0)$ & $\Delta V$ & $\Delta V / V(0)$ \\
    \midrule
    1\,week  & \tbd{2.14} & \tbd{2.21} & \tbd{$-$0.07} & \tbd{$-$3.2\%} \\
    1\,month & \tbd{4.83} & \tbd{5.02} & \tbd{$-$0.19} & \tbd{$-$3.8\%} \\
    3\,months& \tbd{8.71} & \tbd{9.14} & \tbd{$-$0.43} & \tbd{$-$4.7\%} \\
    6\,months& \tbd{11.2} & \tbd{11.9} & \tbd{$-$0.70} & \tbd{$-$5.9\%} \\
    \bottomrule
  \end{tabular}
\end{table}

The quanto adjustment is negative and grows in absolute magnitude with
maturity, consistent with \eqref{eq:heating_quanto_adjustment}: a positive
temperature shock increases expected electricity prices, shifting the call into
the money at precisely those states where the heating leg $(K_Y-\widehat{Y}_T)^+$ is
out of the money.
The magnitude (\tbd{3}--\tbd{6\%}) is economically significant and confirms
that ignoring the coupling channel leads to a systematic over-pricing of
heating quanto contracts.

\subsection{Hedging performance}\label{ssec:hedging}

We assess the variance-optimal hedge of \Cref{sec:hedging-strategies} in an
out-of-sample backtest on the 2025 test year.
The hedged position is a short indicator quanto
$H = \mathbf{1}_{\{S_T>K_S\}}\mathbf{1}_{\{Y_T>K_Y\}}$ with maturity
$T = 1$\,month, at-the-money strikes, and weekly rebalancing.
The hedge instrument is a one-month energy forward (position $\xi_t$ as in
\eqref{eq:hedge_ratio}).

Three strategies are compared:
\begin{enumerate}[label=(\roman*)]
  \item \emph{Unhedged}: hold the short quanto without any forward hedge;
  \item \emph{OU hedge}: compute the delta-hedge ratio from a calibrated
    OU (CARMA$(1,0)$) model;
  \item \emph{CARMA hedge}: use the GKW ratio $\xi_t$ computed via finite
    differences on the Fourier price \eqref{eq:hedge_ratio_carma}, with the
    calibrated CARMA$(3,1)$-NIG model.
\end{enumerate}

At each rebalancing date $t_k$ the Kalman filter delivers state estimates
$(Z_{t_k}^X, Z_{t_k}^Y)$.
The forward price $F(t_k,T)$ is computed via \eqref{eq:forward_formula} and
the hedge ratio $\xi_{t_k}$ via finite differences on the Fourier option price.
The weekly hedge PnL from date $t_k$ to $t_{k+1}$ is
$$
  \Pi_{k+1}
  = V_{t_{k+1}} - V_{t_k} - \xi_{t_k}(F_{t_{k+1}} - F_{t_k}).
$$
Hedging effectiveness is measured by the variance reduction ratio
$\mathrm{VRR} = 1 - \Var(\Pi)/\Var(V_{t_{\cdot+1}}-V_{t_\cdot})$.

\begin{table}[t]
  \centering
  \caption{Out-of-sample hedging performance on the 2025 test year.
           Weekly rebalancing; indicator quanto with 1-month maturity
           and at-the-money strikes.
           VRR $= 1 - \Var(\Pi)/\Var(\Delta V)$; RMSE is the root-mean-square
           weekly hedge error.}
  \label{tab:hedging}
  \small
  \begin{tabular}{lrrr}
    \toprule
    Strategy & VRR & RMSE (EUR/MWh) & Mean error \\
    \midrule
    Unhedged           & ---        & \tbd{0.0182} & \tbd{$-$0.0003} \\
    OU hedge           & \tbd{28\%} & \tbd{0.0155} & \tbd{$-$0.0002} \\
    CARMA$(3,1)$ hedge & \tbd{41\%} & \tbd{0.0135} & \tbd{$-$0.0001} \\
    \bottomrule
  \end{tabular}
\end{table}

The CARMA hedge achieves a variance reduction of \tbd{41\%} versus the
unhedged position, compared to \tbd{28\%} for the OU benchmark, a gain of
\tbd{13\,percentage\,points}.
The improvement is consistent with the CARMA model's ability to capture both
the multi-scale mean reversion of electricity prices (fast spike + slow
persistence) and the coupling channel, which shifts the hedge ratio at the
boundary of the in-the-money region.
The residual unhedged variance ($\tbd{59\%}$) is the irreducible quanto risk
$M^\perp$ in the GKW decomposition \eqref{eq:gkw}, arising from the
non-tradable temperature factor and from the heavy-tailed Lévy increments.

\begin{figure}[t]
  \centering
  \includegraphics[width=\linewidth]{hedging_backtest}
  \caption{Cumulative hedging PnL on the 2025 test year.
           Solid blue: unhedged short quanto position.
           Dashed orange: OU-hedged position.
           Solid red: CARMA$(3,1)$-NIG hedged position.
           The shaded band is the unhedged $\pm$1\,std.\ dev.\ envelope.}
  \label{fig:hedging}
\end{figure}
